\documentclass[10pt,fleqn]{article} % Default font size and left-justified equations
\usepackage[%
    pdftitle={SLCI : Modélisation par schémas blocs},
    pdfauthor={Xavier Pessoles}]{hyperref}

%%%%%%%%%%%%%%%%%%%%%%%%%%%%%%%%%%%%%%%%%
% Original author:
% Mathias Legrand (legrand.mathias@gmail.com) with modifications by:
% Vel (vel@latextemplates.com)
% License:
% CC BY-NC-SA 3.0 (http://creativecommons.org/licenses/by-nc-sa/3.0/)
%%%%%%%%%%%%%%%%%%%%%%%%%%%%%%%%%%%%%%%%%

%----------------------------------------------------------------------------------------
%	VARIOUS REQUIRED PACKAGES AND CONFIGURATIONS
%----------------------------------------------------------------------------------------

\usepackage[top=2.5cm,bottom=2cm,left=2cm,right=2cm,headsep=40pt,a4paper]{geometry} % Page margins

\usepackage{graphicx} % Required for including pictures
\graphicspath{{images/}} % Specifies the directory where pictures are stored

\usepackage{lipsum} % Inserts dummy text

\usepackage{tikz} % Required for drawing custom shapes

\usepackage[french]{babel} % English language/hyphenation
\frenchbsetup{StandardLists=true} % Pour éviter la collision babel enumitem pour les listes

\usepackage{enumitem} % Customize lists
\setlist{nolistsep} % Reduce spacing between bullet points and numbered lists

\usepackage{booktabs} % Required for nicer horizontal rules in tables

\usepackage{xcolor} % Required for specifying colors by name
%\definecolor{ocre}{RGB}{243,102,25} % Define the orange color used for highlighting throughout the book
 \definecolor{ocre}{RGB}{49,133,156} % Couleur ''bleue''
\definecolor{violetf}{RGB}{112,48,160} % Couleur ''violet''
\usepackage{enumitem}
\usepackage{pifont} % Pour les dinglist
\usepackage{multicol}
\usepackage{array} % Centrage vertical dans les tableaux

%----------------------------------------------------------------------------------------
%	FONTS
%----------------------------------------------------------------------------------------

\usepackage{avant} % Use the Avantgarde font for headings
%\usepackage{times} % Use the Times font for headings
%\usepackage{mathptmx} % Use the Adobe Times Roman as the default text font together with math symbols from the Sym­bol, Chancery and Com­puter Modern fonts
\usepackage[adobe-utopia]{mathdesign}
\usepackage{microtype} % Slightly tweak font spacing for aesthetics
\usepackage[utf8]{inputenc} % Required for including letters with accents
\usepackage[T1]{fontenc} % Use 8-bit encoding that has 256 glyphs

%----------------------------------------------------------------------------------------
%	BIBLIOGRAPHY AND INDEX
%----------------------------------------------------------------------------------------

\usepackage[style=alphabetic,citestyle=numeric,sorting=nyt,sortcites=true,autopunct=true,babel=hyphen,hyperref=true,abbreviate=false,backref=true,backend=biber]{biblatex}
\addbibresource{bibliography.bib} % BibTeX bibliography file
\defbibheading{bibempty}{}

\usepackage{calc} % For simpler calculation - used for spacing the index letter headings correctly
\usepackage{makeidx} % Required to make an index
\makeindex % Tells LaTeX to create the files required for indexing

%----------------------------------------------------------------------------------------
%	MAIN TABLE OF CONTENTS
%----------------------------------------------------------------------------------------

\usepackage{titletoc} % Required for manipulating the table of contents

\setcounter{tocdepth}{2}     % Dans la table des matieres
\setcounter{secnumdepth}{2}

\contentsmargin{0cm} % Removes the default margin

% Part text styling
\titlecontents{part}[0cm]
{\addvspace{20pt}\centering\large\bfseries}
{}
{}
{}

% Chapter text styling
\titlecontents{chapter}[1.25cm] % Indentation
{\addvspace{12pt}\large\sffamily\bfseries} % Spacing and font options for chapters
{\color{ocre!60}\contentslabel[\Large\thecontentslabel]{1.25cm}\color{ocre}} % Chapter number
{\color{ocre}}  
{\color{ocre!60}\normalsize\;\titlerule*[.5pc]{.}\;\thecontentspage} % Page number

% Section text styling
\titlecontents{section}[1.25cm] % Indentation
{\addvspace{3pt}\sffamily\bfseries} % Spacing and font options for sections
{\color{ocre!60}\contentslabel[\thecontentslabel]{1.25cm} \color{ocre}} % Section number
{\color{ocre}}
{\hfill\color{ocre!60}\thecontentspage} % Page number
[]

% Subsection text styling
\titlecontents{subsection}[1.25cm] % Indentation
{\addvspace{1pt}\sffamily\small} % Spacing and font options for subsections
{\contentslabel[\thecontentslabel]{1.25cm}} % Subsection number
{}
{\ \titlerule*[.5pc]{.}\;\thecontentspage} % Page number
[]


% Subsection text styling
\titlecontents{subsubsection}[1.25cm] % Indentation
{\addvspace{1pt}\sffamily\small} % Spacing and font options for subsections
{\contentslabel[\thecontentslabel]{1.25cm}} % Subsection number
{}
{\ \titlerule*[.5pc]{.}\;\thecontentspage} % Page number
[]

% List of figures
\titlecontents{figure}[0em]
{\addvspace{-5pt}\sffamily}
{\thecontentslabel\hspace*{1em}}
{}
{\ \titlerule*[.5pc]{.}\;\thecontentspage}
[]

% List of tables
\titlecontents{table}[0em]
{\addvspace{-5pt}\sffamily}
{\thecontentslabel\hspace*{1em}}
{}
{\ \titlerule*[.5pc]{.}\;\thecontentspage}
[]

%----------------------------------------------------------------------------------------
%	MINI TABLE OF CONTENTS IN PART HEADS
%----------------------------------------------------------------------------------------

% Chapter text styling
\titlecontents{lchapter}[0em] % Indenting
{\addvspace{15pt}\large\sffamily\bfseries} % Spacing and font options for chapters
{\color{ocre}\contentslabel[\Large\thecontentslabel]{1.25cm}\color{ocre}} % Chapter number
{}  
{\color{ocre}\normalsize\sffamily\bfseries\;\titlerule*[.5pc]{.}\;\thecontentspage} % Page number

% Section text styling
\titlecontents{lsection}[0em] % Indenting
{\sffamily\small} % Spacing and font options for sections
{\contentslabel[\thecontentslabel]{1.25cm}} % Section number
{}
{}

% Subsection text styling
\titlecontents{lsubsection}[.5em] % Indentation
{\normalfont\footnotesize\sffamily} % Font settings
{}
{}
{}

%----------------------------------------------------------------------------------------
%	PAGE HEADERS
%----------------------------------------------------------------------------------------

\usepackage{fancyhdr} % Required for header and footer configuration



\pagestyle{fancy}
 \renewcommand{\headrulewidth}{0pt}
 \fancyhead{}
 \fancyhead[L]{%
 \noindent\begin{minipage}[c]{2.6cm}%
 \includegraphics[width=2cm]{png/logo_lycee.png}%
 \end{minipage}}

\fancyhead[C]{\rule{8cm}{.5pt}}

 \fancyhead[R]{%
 \noindent\begin{minipage}[c]{3cm}
 \begin{flushright}
 \footnotesize{\textit{\textsf{\xxtete}}}%
 \end{flushright}
 \end{minipage}
}


\fancyfoot[C]{\rule{12cm}{.5pt}}
\renewcommand{\footrulewidth}{0.2pt}
\fancyfoot[C]{\footnotesize{\bfseries \thepage}}
\fancyfoot[L]{ 
\begin{minipage}[c]{.4\linewidth}
\noindent\footnotesize{{\xxauteur}}
\end{minipage}}


\fancyfoot[R]{\footnotesize{\xxpied}
\ifthenelse{\isodd{\value{page}}}{
\begin{tikzpicture}[overlay]
\node[shape=rectangle, 
      rounded corners = .25 cm,
	  draw= ocre,
	  line width=2pt, 
	  fill = ocre!10,
	  minimum width  = 2.5cm,
	  minimum height = 3cm,] at (\xxposongletx,\xxposonglety) {};
\node at (\xxposonglettext,\xxposonglety) {\rotatebox{90}{\textbf{\large\color{ocre}{\xxonglet}}}};
%{};
\end{tikzpicture}}{}
}
%
%
%
% Removes the header from odd empty pages at the end of chapters
\makeatletter
\renewcommand{\cleardoublepage}{
\clearpage\ifodd\c@page\else
\hbox{}
\vspace*{\fill}
\thispagestyle{empty}
\newpage
\fi}

\fancypagestyle{plain}{%
\fancyhf{} % vide l’en-tête et le pied~de~page.
%\fancyfoot[C]{\bfseries \thepage} % numéro de la page en cours en gras
% et centré en pied~de~page.
\fancyfoot[R]{\footnotesize{\xxpied}}
\fancyfoot[C]{\rule{12cm}{.5pt}}
\renewcommand{\footrulewidth}{0.2pt}
\fancyfoot[C]{\footnotesize{\bfseries \thepage}}
\fancyfoot[L]{ 
\begin{minipage}[c]{.4\linewidth}
\noindent\footnotesize{{\xxauteur}}
\end{minipage}}}



%----------------------------------------------------------------------------------------
%	THEOREM STYLES
%----------------------------------------------------------------------------------------

% Conflit avec la police adobe
%\usepackage{amsmath,amsfonts,amssymb,amsthm} % For math equations, theorems, symbols, etc
\usepackage{amsmath,amsthm}

\newcommand{\intoo}[2]{\mathopen{]}#1\,;#2\mathclose{[}}
\newcommand{\ud}{\mathop{\mathrm{{}d}}\mathopen{}}
\newcommand{\intff}[2]{\mathopen{[}#1\,;#2\mathclose{]}}
%\newtheorem{notation}{Notation}[chapter]
\newtheorem{notation}{Notation}[section]

% Boxed/framed environments
\newtheoremstyle{ocrenumbox}% % Theorem style name
{0pt}% Space above
{0pt}% Space below
{\normalfont}% % Body font
{}% Indent amount
{\small\bf\sffamily\color{ocre}}% % Theorem head font
{\;}% Punctuation after theorem head
{0.25em}% Space after theorem head
{\small\sffamily\color{ocre}\thmname{#1}\nobreakspace\thmnumber%{\@ifnotempty{#1}{}\@upn{#2}}% Theorem text (e.g. Theorem 2.1)
\thmnote{\nobreakspace\the\thm@notefont\sffamily\bfseries\color{black}---\nobreakspace#3.}} % Optional theorem note
\renewcommand{\qedsymbol}{$\blacksquare$}% Optional qed square


% Boite pour les corriges
\newtheoremstyle{correctionbox}% % Theorem style name
{0pt}% Space above
{0pt}% Space below
{\normalfont}% % Body font
{}% Indent amount
{\small\bf\sffamily\color{violet}}% % Theorem head font
{\;}% Punctuation after theorem head
{0.25em}% Space after theorem head
{\small\sffamily\color{ocre}\thmname{#1}\nobreakspace\thmnumber%{\@ifnotempty{#1}{}\@upn{#2}}% Theorem text (e.g. Theorem 2.1)
\thmnote{\nobreakspace\the\thm@notefont\sffamily\bfseries\color{black}---\nobreakspace#3.}} % Optional theorem note
\renewcommand{\qedsymbol}{$\blacksquare$}% Optional qed square



\newtheoremstyle{blacknumex}% Theorem style name
{5pt}% Space above
{5pt}% Space below
{\normalfont}% Body font
{} % Indent amount
{\small\bf\sffamily}% Theorem head font
{\;}% Punctuation after theorem head
{0.25em}% Space after theorem head
{\small\sffamily{\tiny\ensuremath{\blacksquare}}\nobreakspace\thmname{#1}\nobreakspace\thmnumber%{\@ifnotempty{#1}{}\@upn{#2}}% Theorem text (e.g. Theorem 2.1)
\thmnote{\nobreakspace\the\thm@notefont\sffamily\bfseries---\nobreakspace#3.}}% Optional theorem note

\newtheoremstyle{blacknumbox} % Theorem style name
{0pt}% Space above
{0pt}% Space below
{\normalfont}% Body font
{}% Indent amount
{\small\bf\sffamily}% Theorem head font
{\;}% Punctuation after theorem head
{0.25em}% Space after theorem head
{\small\sffamily\thmname{#1}\nobreakspace 
\thmnote{\nobreakspace\the\thm@notefont\sffamily\bfseries---\nobreakspace#3.}}% Optional theorem note

% Non-boxed/non-framed environments
\newtheoremstyle{ocrenum}% % Theorem style name
{5pt}% Space above
{5pt}% Space below
{\normalfont}% % Body font
{}% Indent amount
{\small\bf\sffamily\color{ocre}}% % Theorem head font
{\;}% Punctuation after theorem head
{0.25em}% Space after theorem head
{\small\sffamily\color{ocre}\thmname{#1}\nobreakspace%\thmnumber{\@ifnotempty{#1}{}\@upn{#2}}% Theorem text (e.g. Theorem 2.1)
\thmnote{\nobreakspace\the\thm@notefont\sffamily\bfseries\color{black}---\nobreakspace#3.}} % Optional theorem note
\renewcommand{\qedsymbol}{$\blacksquare$}% Optional qed square
\makeatother

% Environnement pour les titres de parties
\newtheoremstyle{partiebox} 
{0pt}% Space above
{0pt}% Space below
{\normalfont}% Body font
{}% Indent amount
{\small\bf\sffamily}% Theorem head font
{\;}% Punctuation after theorem head
{0.25em}% Space after theorem head




% Defines the theorem text style for each type of theorem to one of the three styles above
\newcounter{dummy} 
\numberwithin{dummy}{section}
\theoremstyle{ocrenumbox}
%\newtheorem{theoremeT}[dummy]{Théorème}
\newtheorem{theoremeT}[dummy]{Théorème}
\newtheorem{resultatT}[dummy]{Résultat}
\newtheorem{savoirT}[dummy]{Savoir}
\newtheorem{methodeT}[dummy]{Méthode}
\newtheorem{objectifT}[dummy]{Objectif}
%\newtheorem{problem}{Problem}[chapter]
\newtheorem{problem}{Problem}[section]
%\newtheorem{exerciseT}{Exercise}[chapter]
\newtheorem{exerciseT}{Exercice}[section]

\theoremstyle{blacknumex}
%\newtheorem{exampleT}{Example}[chapter]
\newtheorem{exempleT}{Exemple}[section]
\newtheorem{termT}{Terminal\\}[section]
\newtheorem{pyT}{Python\\}[section]
\newtheorem{sciT}{Scilab\\}[section]
\newtheorem{pseudoT}{Pseudo Code\\}[section]
\newtheorem{sqlT}{SQL\\}[section]

\theoremstyle{blacknumbox}
%\newtheorem{vocabulary}{Vocabulary}[chapter]
\newtheorem{vocabulary}{Vocabulaire}[section]
%\newtheorem{definitionT}{Definition}[section]
\newtheorem{definitionT}{Définition}[section]
\newtheorem{rappelT}{Rappel}[section]
\newtheorem{demoT}{Démonstration}[section]
\newtheorem{corollaryT}[dummy]{Corollaire}
\newtheorem{hypoT}{Hypothèse(s)}

\theoremstyle{ocrenum}
\newtheorem{proposition}[dummy]{Proposition}

\theoremstyle{partiebox}
\newtheorem{titrepartieT}[]{}
\newtheorem{titrechapitreT}[]{}

\theoremstyle{correctionbox}
\newtheorem{correctionT}[dummy]{\color{violet}{Correction}}

%----------------------------------------------------------------------------------------
%	DEFINITION OF COLORED BOXES
%----------------------------------------------------------------------------------------

\RequirePackage[framemethod=tikz]{mdframed} % Required for creating the theorem, definition, exercise and corollary boxes

% Theorem box
\newmdenv[skipabove=7pt,
skipbelow=7pt,
backgroundcolor=ocre!10,
linecolor=ocre,
innerleftmargin=5pt,
innerrightmargin=5pt,
innertopmargin=5pt,
leftmargin=0cm,
rightmargin=0cm,
innerbottommargin=5pt]{tBox}


% Correction
\newmdenv[skipabove=7pt,
skipbelow=7pt,
backgroundcolor=violet!10,
linecolor=violet,
innerleftmargin=5pt,
innerrightmargin=5pt,
innertopmargin=5pt,
leftmargin=0cm,
rightmargin=0cm,
innerbottommargin=5pt]{coBox}


% Exercise box	  
\newmdenv[skipabove=7pt,
skipbelow=7pt,
rightline=false,
leftline=true,
topline=false,
bottomline=false,
backgroundcolor=ocre!10,
linecolor=ocre,
innerleftmargin=5pt,
innerrightmargin=5pt,
innertopmargin=5pt,
innerbottommargin=5pt,
leftmargin=0cm,
rightmargin=0cm,
linewidth=4pt]{eBox}	

% Definition box
\newmdenv[skipabove=7pt,
skipbelow=7pt,
rightline=false,
leftline=true,
topline=false,
bottomline=false,
backgroundcolor=ocre!10,
linecolor=ocre,
innerleftmargin=5pt,
innerrightmargin=5pt,
innertopmargin=0pt,
leftmargin=0cm,
rightmargin=0cm,
linewidth=4pt,
innerbottommargin=0pt]{dBox}	

% Demonstration box
\newmdenv[skipabove=7pt,
skipbelow=7pt,
rightline=false,
leftline=true,
topline=false,
bottomline=false,
%backgroundcolor=ocre!10,
linecolor=ocre,
innerleftmargin=5pt,
innerrightmargin=5pt,
innertopmargin=0pt,
leftmargin=0cm,
rightmargin=0cm,
linewidth=4pt,
innerbottommargin=0pt]{demoBox}	

% Corollary box
\newmdenv[skipabove=7pt,
skipbelow=7pt,
rightline=false,
leftline=true,
topline=false,
bottomline=false,
linecolor=gray,
backgroundcolor=black!5,
innerleftmargin=5pt,
innerrightmargin=5pt,
innertopmargin=5pt,
leftmargin=0cm,
rightmargin=0cm,
linewidth=4pt,
innerbottommargin=5pt]{cBox}


% Hypothèses
\newmdenv[skipabove=7pt,
skipbelow=7pt,
rightline=false,
leftline=true,
topline=false,
bottomline=false,
linecolor=gray,
backgroundcolor=black!5,
innerleftmargin=5pt,
innerrightmargin=5pt,
innertopmargin=5pt,
leftmargin=0cm,
rightmargin=0cm,
linewidth=4pt,
innerbottommargin=5pt]{hyBox}


% Boite pour le titre de la partie (pBox)
\newmdenv[skipabove=7pt,
skipbelow=7pt,
rightline=true,
leftline=false,
topline=false,
bottomline=false,
linecolor=ocre,
backgroundcolor=none,
innerleftmargin=5pt,
innerrightmargin=5pt,
innertopmargin=5pt,
leftmargin=0cm,
rightmargin=0cm,
linewidth=4pt,
innerbottommargin=5pt]{pBox}

% Boite pour le titre du chapitre (chBox)
\newmdenv[skipabove=7pt,
skipbelow=7pt,
rightline=false,
leftline=true,
topline=false,
bottomline=false,
linecolor=ocre,
%backgroundcolor=black!5,
innerleftmargin=5pt,
innerrightmargin=5pt,
innertopmargin=5pt,
leftmargin=0cm,
rightmargin=0cm,
linewidth=4pt,
innerbottommargin=5pt]{chBox}


% Boite pour les exemples
\newmdenv[skipabove=7pt,
skipbelow=7pt,
rightline=false,
leftline=true,
topline=false,
bottomline=false,
linecolor=gray,
backgroundcolor=white,
innerleftmargin=5pt,
innerrightmargin=5pt,
innertopmargin=5pt,
leftmargin=0cm,
rightmargin=0cm,
linewidth=4pt,
innerbottommargin=5pt]{exBox}

% Boite pour le terminal
\newmdenv[skipabove=7pt,
skipbelow=7pt,
rightline=false,
leftline=true,
topline=false,
bottomline=false,
linecolor=gray,
backgroundcolor=white,
innerleftmargin=5pt,
innerrightmargin=5pt,
innertopmargin=5pt,
leftmargin=0cm,
rightmargin=0cm,
linewidth=4pt,
innerbottommargin=5pt]{termBox}


% Boite pour Python
\newmdenv[skipabove=7pt,
skipbelow=7pt,
rightline=false,
leftline=true,
topline=false,
bottomline=false,
linecolor=gray,
backgroundcolor=white,
innerleftmargin=5pt,
innerrightmargin=5pt,
innertopmargin=0pt,
leftmargin=0cm,
rightmargin=0cm,
linewidth=4pt,
innerbottommargin=5pt]{pyBox}

% Boite pour scilab
\newmdenv[skipabove=7pt,
skipbelow=7pt,
rightline=false,
leftline=true,
topline=false,
bottomline=false,
linecolor=gray,
backgroundcolor=white,
innerleftmargin=5pt,
innerrightmargin=5pt,
innertopmargin=5pt,
leftmargin=0cm,
rightmargin=0cm,
linewidth=4pt,
innerbottommargin=5pt]{sciBox}


% Boite pour pseudo
\newmdenv[skipabove=7pt,
skipbelow=7pt,
rightline=false,
leftline=true,
topline=false,
bottomline=false,
linecolor=gray,
backgroundcolor=white,
innerleftmargin=5pt,
innerrightmargin=5pt,
innertopmargin=5pt,
leftmargin=0cm,
rightmargin=0cm,
linewidth=4pt,
innerbottommargin=5pt]{pseudoBox}

% Boite pour pseudo
\newmdenv[skipabove=7pt,
skipbelow=7pt,
rightline=false,
leftline=true,
topline=false,
bottomline=false,
linecolor=gray,
backgroundcolor=white,
innerleftmargin=5pt,
innerrightmargin=5pt,
innertopmargin=5pt,
leftmargin=0cm,
rightmargin=0cm,
linewidth=4pt,
innerbottommargin=5pt]{sqlBox}


% Creates an environment for each type of theorem and assigns it a theorem text style from the "Theorem Styles" section above and a colored box from above
\newenvironment{theorem}{\begin{tBox}\begin{theoremeT}}{\end{theoremeT}\end{tBox}}
\newenvironment{resultat}{\begin{tBox}\begin{resultatT}}{\end{resultatT}\end{tBox}}
\newenvironment{methode}{\begin{tBox}\begin{methodeT}}{\end{methodeT}\end{tBox}}
\newenvironment{savoir}{\begin{tBox}\begin{savoirT}}{\end{savoirT}\end{tBox}}
\newenvironment{obj}{\begin{tBox}\begin{objectifT}}{\end{objectifT}\end{tBox}}
\newenvironment{corrige}{\begin{coBox}\begin{correctionT}}{\end{correctionT}\end{coBox}}
\newenvironment{exercise}{\begin{eBox}\begin{exerciseT}}{\hfill{\color{ocre}\tiny\ensuremath{\blacksquare}}\end{exerciseT}\end{eBox}}				  
\newenvironment{exercice}{\begin{eBox}\begin{exerciseT}}{\hfill{\color{ocre}\tiny\ensuremath{\blacksquare}}\end{exerciseT}\end{eBox}}				  

\newenvironment{definition}{\begin{dBox}\begin{definitionT}}{\end{definitionT}\end{dBox}}	
\newenvironment{rappel}{\begin{dBox}\begin{rappelT}}{\end{rappelT}\end{dBox}}	
\newenvironment{defi}{\begin{dBox}\begin{definitionT}}{\end{definitionT}\end{dBox}}	
\newenvironment{demo}{\begin{demoBox}\begin{demoT}}{\end{demoT}\end{demoBox}}	
%\newenvironment{exemple}{\begin{exempleT}}{\hfill{\tiny\ensuremath{\blacksquare}}\end{exempleT}}		
\newenvironment{corollary}{\begin{cBox}\begin{corollaryT}}{\end{corollaryT}\end{cBox}}
\newenvironment{hypo}{\begin{hyBox}\begin{hypoT}}{\end{hypoT}\end{hyBox}}	\newenvironment{exemple}{\begin{exBox}\begin{exempleT}}{\hfill{\tiny\ensuremath{\blacksquare}}\end{exempleT}\end{exBox}}	
\newenvironment{titrepartie}{\begin{pBox}\begin{titrepartieT}}{\end{titrepartieT}\end{pBox}}	
\newenvironment{titrechapitre}{\begin{chBox}\begin{titrechapitreT}}{\end{titrechapitreT}\end{chBox}}	

\newenvironment{term}{ \begin{termBox}\begin{termT}}{\end{termT}\end{termBox}}
\newenvironment{py}{ \begin{pyBox}\begin{pyT}}{\end{pyT}\end{pyBox}}
\newenvironment{sci}{ \begin{sciBox}\begin{sciT}}{\end{sciT}\end{sciBox}}
\newenvironment{pseudo}{ \begin{pseudoBox}\begin{pseudoT}}{\end{pseudoT}\end{pseudoBox}}
\newenvironment{envsql}{ \begin{sqlBox}\begin{sqlT}}{\end{sqlT}\end{sqlBox}}


%----------------------------------------------------------------------------------------
%	REMARK ENVIRONMENT
%----------------------------------------------------------------------------------------

\newenvironment{remark}{\par\vspace{10pt}\small % Vertical white space above the remark and smaller font size
\begin{list}{}{
\leftmargin=35pt % Indentation on the left
\rightmargin=25pt}\item\ignorespaces % Indentation on the right
\makebox[-2.5pt]{\begin{tikzpicture}[overlay]
\node[draw=ocre!60,line width=1pt,circle,fill=ocre!25,font=\sffamily\bfseries,inner sep=2pt,outer sep=0pt] at (-15pt,0pt){\textcolor{ocre}{R}};\end{tikzpicture}} % Orange R in a circle
\advance\baselineskip -1pt}{\end{list}\vskip5pt} % Tighter line spacing and white space after remark

\newenvironment{rem}{\par\vspace{10pt}\small % Vertical white space above the remark and smaller font size
\begin{list}{}{
\leftmargin=35pt % Indentation on the left
\rightmargin=25pt}\item\ignorespaces % Indentation on the right
\makebox[-2.5pt]{\begin{tikzpicture}[overlay]
\node[draw=ocre!60,line width=1pt,circle,fill=ocre!25,font=\sffamily\bfseries,inner sep=2pt,outer sep=0pt] at (-15pt,0pt){\textcolor{ocre}{R}};\end{tikzpicture}} % Orange R in a circle
\advance\baselineskip -1pt}{\end{list}\vskip5pt} % Tighter line spacing and white space after remark


\newenvironment{warn}{\par\vspace{10pt}\small % Vertical white space above the remark and smaller font size
\begin{list}{}{
\leftmargin=35pt % Indentation on the left
\rightmargin=25pt}\item\ignorespaces % Indentation on the right
\makebox[-2.5pt]{\begin{tikzpicture}[overlay]
\node[draw=red!60,line width=1pt,circle,fill=red!25,font=\sffamily\bfseries,inner sep=2pt,outer sep=0pt] at (-15pt,0pt){\textcolor{black}{!}};\end{tikzpicture}} % Point d'exclamation dans un cercle
\advance\baselineskip -1pt}{\end{list}\vskip5pt} % Tighter line spacing and white space after remark


%----------------------------------------------------------------------------------------
%	SECTION NUMBERING IN THE MARGIN
%----------------------------------------------------------------------------------------
\setcounter{secnumdepth}{3}
\setcounter{tocdepth}{2}



\makeatletter
\renewcommand{\@seccntformat}[1]{\llap{\textcolor{ocre}{\csname the#1\endcsname}\hspace{1em}}}                    
\renewcommand{\section}{\@startsection{section}{1}{\z@}
{-4ex \@plus -1ex \@minus -.4ex}
{1ex \@plus.2ex }
{\normalfont\large\sffamily\bfseries}}
\renewcommand{\subsection}{\@startsection {subsection}{2}{\z@}
{-3ex \@plus -0.1ex \@minus -.4ex}
{0.5ex \@plus.2ex }
{\normalfont\sffamily\bfseries}}
\renewcommand{\subsubsection}{\@startsection {subsubsection}{3}{\z@}
{-2ex \@plus -0.1ex \@minus -.2ex}
{.2ex \@plus.2ex }
{\normalfont\small\sffamily\bfseries}}                        
\renewcommand\paragraph{\@startsection{paragraph}{4}{\z@}
{-2ex \@plus-.2ex \@minus .2ex}
{.1ex}
{\normalfont\small\sffamily\bfseries}}

%----------------------------------------------------------------------------------------
%	PART HEADINGS
%----------------------------------------------------------------------------------------


%----------------------------------------------------------------------------------------
%	CHAPTER HEADINGS
%----------------------------------------------------------------------------------------

% \newcommand{\thechapterimage}{}%
% \newcommand{\chapterimage}[1]{\renewcommand{\thechapterimage}{#1}}%
% \def\@makechapterhead#1{%
% {\parindent \z@ \raggedright \normalfont
% \ifnum \c@secnumdepth >\m@ne
% \if@mainmatter
% \begin{tikzpicture}[remember picture,overlay]
% \node at (current page.north west)
% {\begin{tikzpicture}[remember picture,overlay]
% \node[anchor=north west,inner sep=0pt] at (0,0) {\includegraphics[width=\paperwidth]{\thechapterimage}};
% \draw[anchor=west] (\Gm@lmargin,-9cm) node [line width=2pt,rounded corners=15pt,draw=ocre,fill=white,fill opacity=0.5,inner sep=15pt]{\strut\makebox[22cm]{}};
% \draw[anchor=west] (\Gm@lmargin+.3cm,-9cm) node {\huge\sffamily\bfseries\color{black}\thechapter. #1\strut};
% \end{tikzpicture}};
% \end{tikzpicture}
% \else
% \begin{tikzpicture}[remember picture,overlay]
% \node at (current page.north west)
% {\begin{tikzpicture}[remember picture,overlay]
% \node[anchor=north west,inner sep=0pt] at (0,0) {\includegraphics[width=\paperwidth]{\thechapterimage}};
% \draw[anchor=west] (\Gm@lmargin,-9cm) node [line width=2pt,rounded corners=15pt,draw=ocre,fill=white,fill opacity=0.5,inner sep=15pt]{\strut\makebox[22cm]{}};
% \draw[anchor=west] (\Gm@lmargin+.3cm,-9cm) node {\huge\sffamily\bfseries\color{black}#1\strut};
% \end{tikzpicture}};
% \end{tikzpicture}
% \fi\fi\par\vspace*{270\p@}}}

%-------------------------------------------

\def\@makeschapterhead#1{%
\begin{tikzpicture}[remember picture,overlay]
\node at (current page.north west)
{\begin{tikzpicture}[remember picture,overlay]
\node[anchor=north west,inner sep=0pt] at (0,0) {\includegraphics[width=\paperwidth]{\thechapterimage}};
\draw[anchor=west] (\Gm@lmargin,-9cm) node [line width=2pt,rounded corners=15pt,draw=ocre,fill=white,fill opacity=0.5,inner sep=15pt]{\strut\makebox[22cm]{}};
\draw[anchor=west] (\Gm@lmargin+.3cm,-9cm) node {\huge\sffamily\bfseries\color{black}#1\strut};
\end{tikzpicture}};
\end{tikzpicture}
\par\vspace*{270\p@}}
\makeatother

%----------------------------------------------------------------------------------------
%	HYPERLINKS IN THE DOCUMENTS
%----------------------------------------------------------------------------------------


\hypersetup{hidelinks,backref=true,pagebackref=true,hyperindex=true,colorlinks=false,breaklinks=true,urlcolor= ocre,bookmarks=true,bookmarksopen=false,pdftitle={Title},pdfauthor={Author}}
\usepackage{bookmark}
\bookmarksetup{
open,
numbered,
addtohook={%
\ifnum\bookmarkget{level}=0 % chapter
\bookmarksetup{bold}%
\fi
\ifnum\bookmarkget{level}=-1 % part
\bookmarksetup{color=ocre,bold}%
\fi
}
}

%----------------------------------------------------------------------------------------
%	
%----------------------------------------------------------------------------------------

\newcommand{\thechapterimage}{}%
\newcommand{\chapterimage}[1]{\renewcommand{\thechapterimage}{#1}}%
\def\@makechapterhead#1{%
{\parindent \z@ \raggedright \normalfont
\begin{tikzpicture}[remember picture,overlay]
\node at (current page.north west)
{\begin{tikzpicture}[remember picture,overlay]
\node[anchor=north west,inner sep=0pt] at (0,0) {\includegraphics[width=\paperwidth]{\thechapterimage}};
%\draw[anchor=west] (\Gm@lmargin,-9cm) node [line width=2pt,rounded corners=15pt,draw=ocre,fill=white,fill opacity=0.5,inner sep=15pt]{\strut\makebox[22cm]{}};
%\draw[anchor=west] (\Gm@lmargin+.3cm,-9cm) node {\huge\sffamily\bfseries\color{black}\thechapter. #1\strut};
\end{tikzpicture}};
\end{tikzpicture}
\par\vspace*{270\p@}
}}

 \newcounter{exo}


\makeatletter             
\renewcommand{\subparagraph}{\@startsection{exo}{5}{\z@}%
                                    {-2ex \@plus-.2ex \@minus .2ex}%
                                    {0ex}%               
{\normalfont\bfseries Question \hspace{.7cm} }}
\makeatother
\renewcommand{\thesubparagraph}{\arabic{subparagraph}} 
\makeatletter


%%%% Environnement pour inclure du code
\usepackage{textcomp}
\usepackage[french]{algorithm2e}
\usepackage{listings}
\lstloadlanguages{R}   % pour regler les pb d accent utf8 dans les codes
\lstset{language=R} % pour regler les pb d accent utf8 dans les codes
\renewcommand{\lstlistlistingname}{Listings}
\renewcommand{\lstlistingname}{Listing}

\SetKwBlock{Fonction}{Début Fonction}{Fin Fonction}
\SetKwComment{Comment}{start}{end}

\definecolor{Bleu}{rgb}{0.1,0.1,1.0}
\definecolor{Noir}{rgb}{0,0,0}
\definecolor{Grau}{rgb}{0.5,0.5,0.5}
\definecolor{DunkelGrau}{rgb}{0.15,0.15,0.15}
\definecolor{Hellbraun}{rgb}{0.5,0.25,0.0}
\definecolor{Magenta}{rgb}{1.0,0.0,1.0}
\definecolor{Gris}{gray}{0.5}
\definecolor{Vert}{rgb}{0,0.5,0}
\definecolor{SourceHintergrund}{rgb}{1,1.0,0.95}


\lstnewenvironment{python}[1][]{
\lstset{
%escapeinside={\%*}{*)},
inputencoding=utf8,   % pour regler les pb d accent utf8 dans les codes
extendedchars=true,   % pour regler les pb d accent utf8 dans les codes
language=python,
basicstyle=\ttfamily\footnotesize, 	
stringstyle=\color{red}, 
showstringspaces=false, 
alsoletter={1234567890},
otherkeywords={\ , \}, \{},
keywordstyle=\color{blue},
emph={access,and,break,class,continue,def,del,elif ,else,
except,exec,finally,for,from,global,if,import,in,i s,
lambda,not,or,pass,print,raise,return,try,while},
emphstyle=\color{black}\bfseries,
emph={[2]True, False, None, self},
emphstyle=[2]\color{black},
emph={[3]from, import, as},
emphstyle=[3]\color{blue},
upquote=true,
columns=flexible, % pour empecher d'avoir un espacement mono
morecomment=[s]{"""}{"""},
commentstyle=\color{Hellbraun}\slshape, 
%emph={[4]1, 2, 3, 4, 5, 6, 7, 8, 9, 0},
emphstyle=[4]\color{blue},
literate=*{:}{{\textcolor{blue}:}}{1}
{=}{{\textcolor{blue}=}}{1}
{-}{{\textcolor{blue}-}}{1}
{+}{{\textcolor{blue}+}}{1}
{*}{{\textcolor{blue}*}}{1}
{!}{{\textcolor{blue}!}}{1}
{(}{{\textcolor{blue}(}}{1}
{)}{{\textcolor{blue})}}{1}
{[}{{\textcolor{blue}[}}{1}
{]}{{\textcolor{blue}]}}{1}
{<}{{\textcolor{blue}<}}{1}
{>}{{\textcolor{blue}>}}{1}
{COMPLETER}{{\textcolor{red}COMPLETER}}{1},
literate=%
            {é}{{\'{e}}}1
            {è}{{\`{e}}}1
            {ê}{{\^{e}}}1
            {ë}{{\¨{e}}}1
            {û}{{\^{u}}}1
            {ù}{{\`{u}}}1
            {â}{{\^{a}}}1
            {à}{{\`{a}}}1
            {î}{{\^{i}}}1
            {ç}{{\c{c}}}1
            {Ç}{{\c{C}}}1
            {É}{{\'{E}}}1
            {Ê}{{\^{E}}}1
            {À}{{\`{A}}}1
            {Â}{{\^{A}}}1
            {Î}{{\^{I}}}1, % pour regler les pb d accent utf8 dans les codes
%framexleftmargin=1mm, framextopmargin=1mm, frame=shadowbox, rulesepcolor=\color{blue},#1
%backgroundcolor=\color{SourceHintergrund}, 
%framexleftmargin=1mm, framexrightmargin=1mm, framextopmargin=1mm, frame=single, framerule=1pt, rulecolor=\color{black},#1
}}{}



\lstnewenvironment{scilab}[1][]{
\lstset{
language=scilab,
basicstyle=\sffamily\footnotesize, 	
stringstyle=\color{red}, 
showstringspaces=false, 
alsoletter={1234567890},
otherkeywords={\ , \}, \{},
keywordstyle=\color{blue},
emph={access,and,break,class,continue,def,del,elif ,else,
except,exec,finally,for,from,global,if,import,in,i s,
lambda,not,or,pass,print,raise,return,try,while,Debut},
emphstyle=\color{black}\bfseries,
emph={[2]True, False, None, self},
emphstyle=[2]\color{black},
emph={[3]from, import, as},
emphstyle=[3]\color{blue},
upquote=true,
columns=flexible, % pour empecher d'avoir un espacement mono
morecomment=[s]{"""}{"""},
commentstyle=\color{Hellbraun}\slshape, 
%emph={[4]1, 2, 3, 4, 5, 6, 7, 8, 9, 0},
emphstyle=[4]\color{blue},
literate=*{:}{{\textcolor{blue}:}}{1}
{=}{{\textcolor{blue}=}}{1}
{-}{{\textcolor{blue}-}}{1}
{+}{{\textcolor{blue}+}}{1}
{*}{{\textcolor{blue}*}}{1}
{!}{{\textcolor{blue}!}}{1}
{(}{{\textcolor{blue}(}}{1}
{)}{{\textcolor{blue})}}{1}
{[}{{\textcolor{blue}[}}{1}
{]}{{\textcolor{blue}]}}{1}
{<}{{\textcolor{blue}<}}{1}
{>}{{\textcolor{blue}>}}{1},
%framexleftmargin=1mm, framextopmargin=1mm, frame=shadowbox, rulesepcolor=\color{blue},#1
%backgroundcolor=\color{SourceHintergrund}, 
%framexleftmargin=1mm, framexrightmargin=1mm, framextopmargin=1mm, frame=single, framerule=1pt, rulecolor=\color{black},#1
}}{}


\lstdefinestyle{stylepython}{%
escapeinside={\%*}{*)},
inputencoding=utf8,   % pour regler les pb d accent utf8 dans les codes
extendedchars=true,   % pour regler les pb d accent utf8 dans les codes
language=python,
basicstyle=\sffamily\footnotesize, 	
stringstyle=\color{red}, 
showstringspaces=false, 
alsoletter={1234567890},
otherkeywords={\ , \}, \{},
keywordstyle=\color{blue},
emph={access,and,break,class,continue,def,del,elif ,else,
except,exec,finally,for,from,global,if,import,in,i s,
lambda,not,or,pass,print,raise,return,try,while},
emphstyle=\color{black}\bfseries,
emph={[2]True, False, None, self},
emphstyle=[2]\color{green},
emph={[3]from, import, as},
emphstyle=[3]\color{blue},
upquote=true,
columns=flexible, % pour empecher d'avoir un espacement mono
morecomment=[s]{"""}{"""},
commentstyle=\color{Hellbraun}\slshape, 
%emph={[4]1, 2, 3, 4, 5, 6, 7, 8, 9, 0},
emphstyle=[4]\color{blue},
literate=*{:}{{\textcolor{blue}:}}{1}
{=}{{\textcolor{blue}=}}{1}
{-}{{\textcolor{blue}-}}{1}
{+}{{\textcolor{blue}+}}{1}
{*}{{\textcolor{blue}*}}{1}
{!}{{\textcolor{blue}!}}{1}
{(}{{\textcolor{blue}(}}{1}
{)}{{\textcolor{blue})}}{1}
{[}{{\textcolor{blue}[}}{1}
{]}{{\textcolor{blue}]}}{1}
{<}{{\textcolor{blue}<}}{1}
{>}{{\textcolor{blue}>}}{1}
{COMPLETER}{{\textcolor{red}COMPLETER}}{1},
literate=%
            {é}{{\'{e}}}1
            {è}{{\`{e}}}1
            {ê}{{\^{e}}}1
            {ë}{{\¨{e}}}1
            {û}{{\^{u}}}1
            {ù}{{\`{u}}}1
            {â}{{\^{a}}}1
            {à}{{\`{a}}}1
            {î}{{\^{i}}}1
            {ç}{{\c{c}}}1
            {Ç}{{\c{C}}}1
            {É}{{\'{E}}}1
            {Ê}{{\^{E}}}1
            {À}{{\`{A}}}1
            {Â}{{\^{A}}}1
            {Î}{{\^{I}}}1,
%numbers=left,                    % where to put the line-numbers; possible values are (none, left, right)
%numbersep=5pt,                   % how far the line-numbers are from the code
%numberstyle=\tiny\color{mygray}, % the style that is used for the line-numbers
}



\lstnewenvironment{termi}[1][]{
\lstset{
language=scilab,
basicstyle=\sffamily\footnotesize, 	
stringstyle=\color{red}, 
showstringspaces=false, 
alsoletter={1234567890},
otherkeywords={\ , \}, \{},
keywordstyle=\color{blue},
emph={access,and,break,class,continue,def,del,elif ,else,
except,exec,finally,for,from,global,if,import,in,i s,
lambda,not,or,pass,print,raise,return,try,while,Debut},
emphstyle=\color{black}\bfseries,
emph={[2]True, False, None, self},
emphstyle=[2]\color{green},
emph={[3]from, import, as},
emphstyle=[3]\color{blue},
upquote=true,
columns=flexible, % pour empecher d'avoir un espacement mono
morecomment=[s]{"""}{"""},
commentstyle=\color{Hellbraun}\slshape, 
%emph={[4]1, 2, 3, 4, 5, 6, 7, 8, 9, 0},
emphstyle=[4]\color{blue},
literate=*{:}{{\textcolor{blue}:}}{1}
{=}{{\textcolor{blue}=}}{1}
{-}{{\textcolor{blue}-}}{1}
{+}{{\textcolor{blue}+}}{1}
{*}{{\textcolor{blue}*}}{1}
{!}{{\textcolor{blue}!}}{1}
{(}{{\textcolor{blue}(}}{1}
{)}{{\textcolor{blue})}}{1}
{[}{{\textcolor{blue}[}}{1}
{]}{{\textcolor{blue}]}}{1}
{<}{{\textcolor{blue}<}}{1}
{>}{{\textcolor{blue}>}}{1},
%framexleftmargin=1mm, framextopmargin=1mm, frame=shadowbox, rulesepcolor=\color{blue},#1
%backgroundcolor=\color{SourceHintergrund}, 
%framexleftmargin=1mm, framexrightmargin=1mm, framextopmargin=1mm, frame=single, framerule=1pt, rulecolor=\color{black},#1
}}{}


\lstnewenvironment{sql}[1][]{
\lstset{
%escapeinside={\%*}{*)},
%inputencoding=utf8,   % pour regler les pb d accent utf8 dans les codes
%extendedchars=true,   % pour regler les pb d accent utf8 dans les codes
language=sql,
basicstyle=\sffamily\footnotesize, 	
stringstyle=\color{red}, 
showstringspaces=false, 
alsoletter={1234567890},
otherkeywords={\ , \}, \{},
keywordstyle=\color{blue},
emph={access,and,break,class,continue,def,del,elif ,else,
except,exec,finally,for,from,global,if,import,in,i s,
lambda,not,or,pass,print,raise,return,try,while},
emphstyle=\color{black}\bfseries,
emph={[2]True, False, None, self},
emphstyle=[2]\color{black},
emph={[3]from, import, as},
emphstyle=[3]\color{blue},
upquote=true,
columns=flexible, % pour empecher d'avoir un espacement mono
morecomment=[s]{"""}{"""},
commentstyle=\color{Hellbraun}\slshape, 
%emph={[4]1, 2, 3, 4, 5, 6, 7, 8, 9, 0},
emphstyle=[4]\color{blue},
literate=*{:}{{\textcolor{blue}:}}{1}
{=}{{\textcolor{blue}=}}{1}
{-}{{\textcolor{blue}-}}{1}
{+}{{\textcolor{blue}+}}{1}
{*}{{\textcolor{blue}*}}{1}
{!}{{\textcolor{blue}!}}{1}
{(}{{\textcolor{blue}(}}{1}
{)}{{\textcolor{blue})}}{1}
{[}{{\textcolor{blue}[}}{1}
{]}{{\textcolor{blue}]}}{1}
{<}{{\textcolor{blue}<}}{1}
{>}{{\textcolor{blue}>}}{1}
{COMPLETER}{{\textcolor{red}COMPLETER}}{1},
literate=%
            {é}{{\'{e}}}1
            {è}{{\`{e}}}1
            {ê}{{\^{e}}}1
            {ë}{{\¨{e}}}1
            {û}{{\^{u}}}1
            {ù}{{\`{u}}}1
            {â}{{\^{a}}}1
            {à}{{\`{a}}}1
            {î}{{\^{i}}}1
            {ç}{{\c{c}}}1
            {Ç}{{\c{C}}}1
            {É}{{\'{E}}}1
            {Ê}{{\^{E}}}1
            {À}{{\`{A}}}1
            {Â}{{\^{A}}}1
            {Î}{{\^{I}}}1, % pour regler les pb d accent utf8 dans les codes
%framexleftmargin=1mm, framextopmargin=1mm, frame=shadowbox, rulesepcolor=\color{blue},#1
%backgroundcolor=\color{SourceHintergrund}, 
%framexleftmargin=1mm, framexrightmargin=1mm, framextopmargin=1mm, frame=single, framerule=1pt, rulecolor=\color{black},#1
}}{}


% Définition des booleéns
\newif\iffiche
\newif\ifprof
\newif\iftd
\newif\ifcours

%%%%%%%%%%%%
% Définition des vecteurs 
%%%%%%%%%%%%
 \newcommand{\vect}[1]{\overrightarrow{#1}}
\newcommand{\axe}[2]{\left(#1,\vect{#2}\right)}

\newcommand{\rep}[1]{\mathcal{R}_{#1}}
\newcommand{\vx}[1]{\vect{x_{#1}}}
\newcommand{\vy}[1]{\vect{y_{#1}}}
\newcommand{\vz}[1]{\vect{z_{#1}}}

%%%%%%%%%%%%
% Définition des torseurs 
%%%%%%%%%%%%

 \newcommand{\torseur}[1]{%
\left\{{#1}\right\}
}

\newcommand{\torseurcin}[3]{%
\left\{\mathcal{#1} \left(#2/#3 \right) \right\}
}

\newcommand{\torseurstat}[3]{%
\left\{\mathcal{#1} \left(#2\rightarrow #3 \right) \right\}
}

 \newcommand{\torseurc}[8]{%
%\left\{#1 \right\}=
\left\{
{#1}
\right\}
 = 
\left\{%
\begin{array}{cc}%
{#2} & {#5}\\%
{#3} & {#6}\\%
{#4} & {#7}\\%
\end{array}%
\right\}_{#8}%
}

 \newcommand{\torseurcol}[7]{
\left\{%
\begin{array}{cc}%
{#1} & {#4}\\%
{#2} & {#5}\\%
{#3} & {#6}\\%
\end{array}%
\right\}_{#7}%
}

 \newcommand{\torseurl}[3]{%
%\left\{\mathcal{#1}\right\}_{#2}=%
\left\{%
\begin{array}{l}%
{#1} \\%
{#2} %
\end{array}%
\right\}_{#3}%
}

 \newcommand{\vectv}[3]{%
\vect{V\left( {#1} \in {#2}/{#3}\right)}
}


\newcommand{\vectf}[2]{%
\vect{R\left( {#1} \rightarrow {#2}\right)}
}

\newcommand{\vectm}[3]{%
\vect{\mathcal{M}\left( {#1}, {#2} \rightarrow {#3}\right)}
}


 \newcommand{\vectg}[3]{%
\vect{\Gamma \left( {#1} \in {#2}/{#3}\right)}
}

 \newcommand{\vecto}[2]{%
\vect{\Omega\left( {#1}/{#2}\right)}
}
% }$$\left\{\mathcal{#1} \right\}_{#2} =%
% \left\{%
% \begin{array}{c}%
%  #3 \\%
%  #4 %
% \end{array}%
% \right\}_{#5}}

\usepackage{style/schemabloc}
%\fichetrue
\fichefalse

\proftrue
%\proffalse

%\tdtrue
\tdfalse

\courstrue
%\coursfalse

% -------------------------------------
% Déclaration des titres
% -------------------------------------

\def\discipline{Sciences \\Industrielles de \\ l'Ingénieur}
\def\xxtete{Sciences Industrielles de l'Ingénieur}

\def\classe{PTSI}
\def\xxnumpartie{Partie 2}
\def\xxpartie{Découverte des Systèmes Linéaires Continus et Invariants\\
Analyse, Modélisation, Résolution}

\def\xxnumchapitre{Chapitre 3}
\def\xxchapitre{\hspace{.12cm} Modélisation des SLCI par schémas blocs}

\def\xxposongletx{2}
\def\xxposonglettext{1.45}
\def\xxposonglety{23}
\def\xxonglet{Part. 2 -- Ch. 3}

\def\xxactivite{Cours}
\def\xxauteur{\textsl{Xavier Pessoles}}

\def\xxcompetences{%
\textsl{%
\textbf{Savoirs et compétences :}
\begin{itemize}[label=\ding{112},font=\color{ocre}] 
\item Modéliser :
\begin{itemize}
\item Mod-C4.1 : Représentation par schéma bloc
\item Mod-C4.2 : Fonction de transfert en boucle ouverte et en boucle fermée
\item Mod-C4.3 : Classe d’un système
\item Mod-C4-S1 : Établir le schéma-bloc du système
\item Mod-C4-S2 : Déterminer les fonctions de transfert du système en boucle ouverte et en boucle fermée.
\end{itemize}
\item Résoudre :
\begin{itemize}
\item Rés-C5-S1 : Schéma bloc ou d’une fonction de transfert les grandeurs caractérisant les performances du modèle
\end{itemize}
\end{itemize}
%
}}

\def\xxfigures{
\begin{center}
\includegraphics[width=4cm]{images/mcc.png}

\textit{Moteur à courant continu}

\includegraphics[width=5cm]{images/modele_mcc.png}
\textit{Schématisation} 

\includegraphics[width=5cm]{images/blocs_mcc.png}
\textit{Schéma bloc}

\end{center}
}%figues de la page de garde
\def\xxpied{%
Partie 2 -- Découverte des SLCI\\
Ch. 3 : Modélisation par schémas blocs -- \xxactivite%
}

%---------------------------------------------------------------------------


\begin{document}
\chapterimage{png/Fond_SLCI}
\pagestyle{empty}


%%%%%%%% PAGE DE GARDE COURS
\ifcours
\begin{tikzpicture}[remember picture,overlay]
\node at (current page.north west)
{\begin{tikzpicture}[remember picture,overlay]
\node[anchor=north west,inner sep=0pt] at (0,0) {\includegraphics[width=\paperwidth]{\thechapterimage}};
\draw[anchor=west] (-2cm,-8cm) node [line width=2pt,rounded corners=15pt,draw=ocre,fill=white,fill opacity=0.6,inner sep=40pt]{\strut\makebox[22cm]{}};
\draw[anchor=west] (1cm,-8cm) node {\huge\sffamily\bfseries\color{black} %
\begin{minipage}{1cm}
\rotatebox{90}{\LARGE\sffamily\textsc{\color{ocre}\textbf{\xxnumpartie}}}
\end{minipage} \hfill
\begin{minipage}[c]{14cm}
\begin{titrepartie}
\begin{flushright}
\renewcommand{\baselinestretch}{1.1} 
\Large\sffamily\textsc{\textbf{\xxpartie}}
\renewcommand{\baselinestretch}{1} 
\end{flushright}
\end{titrepartie}
\end{minipage} \hfill
\begin{minipage}[c]{3.5cm}
{\large\sffamily\textsc{\textbf{\color{ocre} \discipline}}}
\end{minipage} 
 };
\end{tikzpicture}};
\end{tikzpicture}


\begin{tikzpicture}[overlay]
\node[shape=rectangle, 
      rounded corners = .25 cm,
	  draw= ocre,
	  line width=2pt, 
	  fill = ocre!10,
	  minimum width  = 2.5cm,
	  minimum height = 3cm,] at (18cm,0) {};
\node at (17.7cm,0) {\rotatebox{90}{\textbf{\Large\color{ocre}{\classe}}}};
%{};
\end{tikzpicture}

\vspace{3.5cm}

\begin{tikzpicture}[remember picture,overlay]
\draw[anchor=west] (-2cm,-6cm) node {\huge\sffamily\bfseries\color{black} %
\begin{minipage}{2cm}
\begin{center}
\LARGE\sffamily\textsc{\color{ocre}\textbf{\xxactivite}}
\end{center}
\end{minipage} \hfill
\begin{minipage}[c]{15cm}
\begin{titrechapitre}
\renewcommand{\baselinestretch}{1.1} 
\Large\sffamily\textsc{\textbf{\xxnumchapitre}}

\Large\sffamily\textsc{\textbf{\xxchapitre}}
\vspace{.5cm}

\renewcommand{\baselinestretch}{1} 
\normalsize\normalfont
\xxcompetences
\end{titrechapitre}
\end{minipage}  };
\end{tikzpicture}
\vfill

\begin{flushright}
\begin{minipage}[c]{.3\linewidth}
\begin{center}
\xxfigures
\end{center}
\end{minipage}\hfill
\begin{minipage}[c]{.6\linewidth}
\startcontents
\printcontents{}{1}{}
\end{minipage}
\end{flushright}

\begin{tikzpicture}[remember picture,overlay]
\draw[anchor=west] (4.5cm,-.7cm) node {
\begin{minipage}[c]{.2\linewidth}
\begin{flushright}
\includegraphics[width=2cm]{png/logoCC}
\end{flushright}
\end{minipage}
\begin{minipage}[c]{.2\linewidth}
\textsl{\xxauteur} \\
\textsl{\classe}
\end{minipage}
 };
\end{tikzpicture}
\newpage
\pagestyle{fancy}

\newpage
\pagestyle{fancy}

\else
\fi


%%%%%%%% PAGE DE GARDE TD
\iftd
%\begin{tikzpicture}[remember picture,overlay]
%\node at (current page.north west)
%{\begin{tikzpicture}[remember picture,overlay]
%\draw[anchor=west] (-2cm,-3.25cm) node [line width=2pt,rounded corners=15pt,draw=ocre,fill=white,fill opacity=0.6,inner sep=40pt]{\strut\makebox[22cm]{}};
%\draw[anchor=west] (1cm,-3.25cm) node {\huge\sffamily\bfseries\color{black} %
%\begin{minipage}{1cm}
%\rotatebox{90}{\LARGE\sffamily\textsc{\color{ocre}\textbf{\xxnumpartie}}}
%\end{minipage} \hfill
%\begin{minipage}[c]{13.5cm}
%\begin{titrepartie}
%\begin{flushright}
%\renewcommand{\baselinestretch}{1.1} 
%\Large\sffamily\textsc{\textbf{\xxpartie}}
%\renewcommand{\baselinestretch}{1} 
%\end{flushright}
%\end{titrepartie}
%\end{minipage} \hfill
%\begin{minipage}[c]{3.5cm}
%{\large\sffamily\textsc{\textbf{\color{ocre} \discipline}}}
%\end{minipage} 
% };
%\end{tikzpicture}};
%\end{tikzpicture}

%%%%%%%%%% PAGE DE GARDE TD %%%%%%%%%%%%%%%
%\begin{tikzpicture}[overlay]
%\node[shape=rectangle, 
%      rounded corners = .25 cm,
%	  draw= ocre,
%	  line width=2pt, 
%	  fill = ocre!10,
%	  minimum width  = 2.5cm,
%	  minimum height = 2.5cm,] at (18.5cm,0) {};
%\node at (17.7cm,0) {\rotatebox{90}{\textbf{\Large\color{ocre}{\classe}}}};
%%{};
%\end{tikzpicture}

% PARTIE ET CHAPITRE
%\begin{tikzpicture}[remember picture,overlay]
%\draw[anchor=west] (-1cm,-2.1cm) node {\large\sffamily\bfseries\color{black} %
%\begin{minipage}[c]{15cm}
%\begin{flushleft}
%\xxnumchapitre \\
%\xxchapitre
%\end{flushleft}
%\end{minipage}  };
%\end{tikzpicture}

% Bandeau titre exo
\begin{tikzpicture}[remember picture,overlay]
\draw[anchor=west] (-2cm,-6cm) node {\huge\sffamily\bfseries\color{black} %
\begin{minipage}{5cm}
\begin{center}
\LARGE\sffamily\color{ocre}\textbf{\textsc{\xxactivite}}

\begin{center}
\xxfigures
\end{center}

\end{center}
\end{minipage} \hfill
\begin{minipage}[c]{12cm}
\begin{titrechapitre}
\renewcommand{\baselinestretch}{1.1} 
\large\sffamily\textbf{\textsc{\xxtitreexo}}

\small\sffamily{\textbf{\textit{\color{black!70}\xxsourceexo}}}
\vspace{.5cm}

\renewcommand{\baselinestretch}{1} 
\normalsize\normalfont
\xxcompetences
\end{titrechapitre}
\end{minipage}  };
\end{tikzpicture}

\else
\fi


%%%%%%%% PAGE DE GARDE FICHE
\iffiche
\begin{tikzpicture}[remember picture,overlay]
\node at (current page.north west)
{\begin{tikzpicture}[remember picture,overlay]
\draw[anchor=west] (-2cm,-3.25cm) node [line width=2pt,rounded corners=15pt,draw=ocre,fill=white,fill opacity=0.6,inner sep=40pt]{\strut\makebox[22cm]{}};
\draw[anchor=west] (1cm,-3.25cm) node {\huge\sffamily\bfseries\color{black} %
\begin{minipage}{1cm}
\rotatebox{90}{\LARGE\sffamily\textsc{\color{ocre}\textbf{\xxnumpartie}}}
\end{minipage} \hfill
\begin{minipage}[c]{14cm}
\begin{titrepartie}
\begin{flushright}
\renewcommand{\baselinestretch}{1.1} 
\large\sffamily\textsc{\textbf{\xxpartie} \\} 

\vspace{.2cm}

\normalsize\sffamily\textsc{\textbf{\xxnumchapitre -- \xxchapitre}}
\renewcommand{\baselinestretch}{1} 
\end{flushright}
\end{titrepartie}
\end{minipage} \hfill
\begin{minipage}[c]{3.5cm}
{\large\sffamily\textsc{\textbf{\color{ocre} \discipline}}}
\end{minipage} 
 };
\end{tikzpicture}};
\end{tikzpicture}


\begin{tikzpicture}[overlay]
\node[shape=rectangle, 
      rounded corners = .25 cm,
	  draw= ocre,
	  line width=2pt, 
	  fill = ocre!10,
	  minimum width  = 2.5cm,
	  minimum height = 2.5cm,] at (18.5cm,0.5cm) {};
%	  minimum height = 2.5cm,] at (18.5cm,0cm) {};
\node at (17.7cm,0.5) {\rotatebox{90}{\textsf{\textbf{\large\color{ocre}{\classe}}}}};
%{};
\end{tikzpicture}



\else
\fi






On a vu précédemment que les équations différentielles régissant le comportement d'un système peuvent être transformées dans le domaine de Laplace dans le but d'être résolues. 

La complexité des systèmes nous poussent à utiliser une représentation schématique : les schémas blocs. Il faudra alors déterminer le lien entre cette représentation et la représentation dans le domaine de Laplace.



%\begin{prob}
%\textbf{Problématiques}
\begin{itemize}
\item Comment modéliser un système complexe multiphysique en utilisant la modélisation en schéma bloc et la modélisation dans le domaine de Laplace ?
\item Comment déterminer la fonction de transfert d'un système dans le but de prévoir son comportement ?
\end{itemize}
%\end{prob}

\section{Fonction de transfert des systèmes}
\subsection{Le moteur à courant continu}

On adopte, pour le moteur à courant continu, la représentation suivante :
\begin{center}
\includegraphics[width=10cm]{images/modele_mcc.png} 
\end{center}

\begin{multicols}{2}
Loi des mailles dans le circuit électrique : 
$$
u(t)=e(t)+R\cdot i(t) + L\dfrac{di(t)}{dt}
$$

Couple de frottement en sortie du moteur : 
$$
c_r(t) = f_v \omega_m(t)
$$



Équation de la dynamique de l'arbre moteur :

$$
J\dfrac{d\omega_m(t)}{dt} = c_m(t) - c_r(t)
$$

Équation électromécanique :
$$
e(t)=K_E \cdot \omega_m(t)
$$

$$
c_m(t)=K_C \cdot i(t)
$$
\end{multicols}

Ainsi, dans le cas du moteur à courant continu, on peut écrire l'équation différentielle liant la tension d'entrée aux bornes du circuit électrique $u(t)$ et la vitesse angulaire $\omega_m(t)$ :

\begin{exemple}
Équation différentielle du MCC :

%\ifthenelse{\boolean{prof}}{%
$$
u(t)= 
 \left( K_e + \dfrac{R f_v}{K_C}\right) \omega_m(t)
+\left( \dfrac{RJ}{K_C} + \dfrac{Lf_v}{K_C} \right)\dfrac{d\omega_m(t)}{dt} 
+\dfrac{LJ}{K_C^2} \dfrac{d^2\omega_m(t)}{dt^2}  
$$

Dans le cas où on néglige l'impact du couple résistant, on a :
$$
u(t)= 
 K_e \omega_m(t)
 +\dfrac{RJ}{K_C} \dfrac{d\omega_m(t)}{dt} 
+\dfrac{LJ}{K_C^2} \dfrac{d^2\omega_m(t)}{dt^2}  
$$
%}{
%\rotatebox{90}{
%\begin{tabular}{p{4cm}}
%\\
%\end{tabular}}
%}
\end{exemple}

\subsection{Définitions}
Soit un système linéaire, continu, invariant et mono-variable d'entrée $e(t)$
et de sortie $s(t)$. Le système est donc modélisable par une équation
différentielle de la forme suivante :
$$
a_0 s(t) + \sum\limits_{i=1}^{n} a_i\dfrac{d^is(t)}{dt^i} = 
b_0 e(t) + \sum\limits_{i=1}^{m} b_i\dfrac{d^ie(t)}{dt^i}
\quad n\geq m
$$

\begin{minipage}[c]{.6\linewidth}
Ce système est modélisable par un schéma bloc : 
\end{minipage}\hfill
\begin{minipage}[c]{.35\linewidth}
\begin{tikzpicture}
\sbEntree{E}
\sbBloc{sys}{$\quad$Système$\quad$}{E} \sbRelier[$e(t)\quad$]{E}{sys}
\sbSortie{S}{sys} \sbRelier[$\quad s(t)$]{sys}{S}
\end{tikzpicture}
\end{minipage}


En se plaçant dans les conditions de Heaviside, l'équation différentielle de
transforme dans le domaine de Laplace :

$$
a_0 S(p) + \sum\limits_{i=1}^{n} a_i p^i S(p) = 
b_0 E(p) + \sum\limits_{i=1}^{m} b_i p^i E(p) 
$$

En factorisant l'expression, on obtient donc : 
$$
S(p) \sum\limits_{i=0}^{n} a_i p^i  = 
E(p) \sum\limits_{i=0}^{m} b_i p^i 
$$


\begin{defi}
 \textbf{Fonctions de transfert -- Transmittance}

On définit la fonction de transfert d'un système la fonction $H(p)$ telle que :
$$
H(p)
=\dfrac{S(p)}{E(p)} 
= \dfrac{\sum\limits_{i=0}^{m} b_i p^i}{\sum\limits_{i=0}^{n} a_i p^i}
=\dfrac{N(p)}{D(p)} 
$$
\end{defi}


\begin{exemple}
Équation du MCC dans le domaine de Laplace :

%\ifthenelse{\boolean{prof}}{%
$$
U(p)= 
 \left( K_e + \dfrac{R f_v}{K_C}\right) \Omega_m(p)
+\left( \dfrac{RJ}{K_C} + \dfrac{Lf_v}{K_C} \right) p \Omega_m(p)
+\dfrac{LJ}{K_C^2}  p^2 \Omega_m(p)
$$

$$ \Longleftrightarrow
U(p)= \Omega_m(p) \left[ 
 \left( K_e + \dfrac{R f_v}{K_C}\right) 
+\left( \dfrac{RJ}{K_C} + \dfrac{Lf_v}{K_C} \right) p 
+\dfrac{LJ}{K_C^2}  p^2 
\right]
$$
%}{
%\rotatebox{90}{
%\begin{tabular}{p{6cm}}
%\\
%\end{tabular}}
%}


Fonction de Transfert du MCC :

$$
H(p)=\dfrac{\Omega_m(p)}{U(p)}
=\dfrac{1}{ \left( K_e + \dfrac{R f_v}{K_C}\right) 
+\left( \dfrac{RJ}{K_C} + \dfrac{Lf_v}{K_C} \right) p 
+\dfrac{LJ}{K_C^2}  p^2 }
=\dfrac{1}{ \left(  \dfrac{K_e K_C +R f_v}{K_C}\right) 
+\left( \dfrac{RJ+Lf_v}{K_C} \right) p 
+\dfrac{LJ}{K_C^2}  p^2 }
$$

Sous sa forme normalisée, la fonction de transfert devient : 
$$
H(p)=\dfrac{\Omega_m(p)}{U(p)}
=\dfrac{\dfrac{K_C}{K_C K_e + R f_v}}{ 1
+\left( \dfrac{RJ + Lf_v}{K_C K_e + R f_v} \right) p 
+\dfrac{LJ}{ K_C^2 K_e + RK_C f_v}  p^2 }
$$

Note : le dénominateur a été factorisé par $K_e + \dfrac{R f_v}{K_C} = \dfrac{K_C K_e + R f_v}{K_C} $.



Dans le cas où on néglige l'impact du couple résistant, on a :
$$
H(p)=\dfrac{\Omega_m(p)}{U(p)}
=\dfrac{
\dfrac{1}{K_e}}{ 1
+ \dfrac{RJ}{K_C K_e } p 
+\dfrac{LJ}{ K_C^2 K_e }  
p^2 }
$$

\end{exemple}

\begin{defi}
 \textbf{Schéma-bloc}

Soit un système d'entrée $E(p)$, de sortie $S(p)$, caractérisé par une fonction
de transfert $H(p)$. Ce système est alors représenté par le schéma bloc suivant
: 
\begin{center}
\begin{tikzpicture}
\sbEntree{E}
\sbBloc{sys}{$\quad$H(p)$\quad$}{E} \sbRelier[$E(p)\quad$]{E}{sys}
\sbSortie{S}{sys} \sbRelier[$\quad S(p)$]{sys}{S}
\end{tikzpicture} 
\end{center}

La relation entrée -- sortie du système se met alors sous la forme : 
$$
S(p) = E(p) \cdot H(p)
$$
\end{defi}


\begin{rem}
Les deux représentations suivantes sont équivalentes :

\begin{minipage}[c]{.45\linewidth}
\begin{center}
\begin{tikzpicture}
\sbEntree{E}
\sbBloc{sys}{$\quad$H(p)$\quad$}{E} \sbRelier[$E(p)\quad$]{E}{sys}
\sbSortie{S}{sys} \sbRelier[$\quad S(p)$]{sys}{S}
\end{tikzpicture} 
\end{center}
\end{minipage}\hfill
\begin{minipage}[c]{.45\linewidth}
\begin{center}
\begin{tikzpicture}
\sbEntree{E}
\sbBloc{sys}{$\quad\dfrac{1}{H(p)}\quad$}{E} \sbRelier[$S(p)\quad$]{E}{sys}
\sbSortie{S}{sys} \sbRelier[$\quad E(p)$]{sys}{S}
\end{tikzpicture} 
\end{center}
\end{minipage}

\end{rem}

\begin{exemple}
Schéma bloc du moteur à courant continu :
\begin{center}
\begin{tikzpicture}
\sbEntree{E}
\sbBloc{sys}{$\quad$H(p)$\quad$}{E} \sbRelier[$U(p)\quad$]{E}{sys}
\sbSortie{S}{sys} \sbRelier[$\quad \Omega_m(p)$]{sys}{S}
\end{tikzpicture} 

On a bien :
$$
\Omega_m(p) = U(p) \cdot H(p)
$$
\end{center}

\end{exemple}





%\begin{props}
$H(p)$ est une fonction rationnelle en $p$. En factorisant le numérateur et le
dénominateur, $H(p)$ peut s'écrire sous cette forme :
$$
H(p) = \dfrac{N(p)}{D(p)} =
K \dfrac{\left(p-z_1 \right)\left(p-z_2 \right)...\left(p-z_m \right)}{
p^{\alpha} \left(p-p_1 \right)\left(p-p_2 \right)...\left(p-p_n \right)}
$$


 \begin{itemize}
 \item Les $z_i$ sont les \textbf{zéros} de la fonction de transfert (réels ou
complexes).
\item Les $p_i$ sont les \textbf{pôles} de la fonction de transfert (réels ou
complexes).
\item \textbf{Le degré de $D(p)$ est appelé ordre $n$ du système ($n\geq m$ pour les
systèmes physiques).}
\item L'équation $D(p)=0$ est appelée équation caractéristique.
\item Le facteur constant $K$ est appelé gain du système.
\item S'il existe une (ou des) racines nulles d'ordre $\alpha$ de $D(p)$, un
terme $p^\alpha$ apparaît au dénominateur. \textbf{$\alpha$ est la classe (ou type) de
la fonction de transfert}. Il correspond au nombre d'intégrations pures du
système.
\end{itemize}
%\end{props}



\subsection{Représentation par schémas blocs}
La représentation par schéma bloc d’un système est déduite de la représentation
par la transformée de Laplace.

Le schéma fonctionnel, ou schéma bloc, est ainsi une représentation graphique du
système d’équations différentielles, ou des relations entre les variables que
décrit le système d’équations différentielles.

%\begin{methode}
\begin{resultat}
Pour établir un schéma fonctionnel la démarche est la suivante :
\begin{enumerate}
\item Appliquer la transformée de Laplace à chaque équation du système différentiel : on obtient alors un système d’équations linéaires
dans le domaine de Laplace, que va traduire le schéma.
\item Rechercher les fonctions de transfert élémentaires et les variables qu’elles relient.
\item Constituer, tracer les schémas en assemblant les blocs des fonctions de transfert élémentaires.
\item Rassembler les schémas, simplifier à l’aide des règles définies un peu plus tard.
\end{enumerate}

%\end{methode}
\end{resultat}

La représentation par le schéma fonctionnel et la fonction de transfert
permettent ainsi de déterminer les
caractéristiques principales du système sans résoudre d'équations
différentielles.

\section{Schémas blocs et fonctions de transfert de systèmes élémentaires}
Pour les équations "simples" de la physique, on peut aisément mettre les théorèmes fondamentaux sous forme de schémas blocs et de fonctions de transfert.
\subsection{Théorèmes de la mécaniques}

\footnotesize{
\begin{center}
%\begin{tabular}{m{3cm}m{3cm}m{3cm}m{3cm}m{3cm}}
\begin{tabular}{p{3cm}p{3cm}p{3cm}p{3cm}p{3cm}}
%\multicolumn{5}{c}{\textbf{Théorèmes de la mécanique}}\\
\begin{center}
Théorème de la résultante dynamique (position)
\end{center}
&
\textit{Exemple : Mouvement d'un solide en translation}%\includegraphics[height=1cm]{images/poids}}
&
$$ F(t)=M\dfrac{d^2x(t)}{dt^2}$$
&
$$ F(p)=Mp^2X(p)$$
&
\begin{center}
\begin{tikzpicture}
\sbEntree{E}
\sbBloc{sys}{$ \quad \dfrac{1}{Mp^2} \quad $}{E} \sbRelier[$ F(p)\quad $]{E}{sys}
\sbSortie{S}{sys} \sbRelier[$ \quad X(p)$]{sys}{S}
\end{tikzpicture}
\end{center} \\
\begin{center}
Théorème de la résultante dynamique (vitesse)
\end{center}
&
%\multirow{2}{*}{%
\begin{center}
\includegraphics[height=1cm]{images/poids}
\end{center}
%}
&
$$  F(t)=M\dfrac{dv(t)}{dt}$$
&
$$F(p)=MpV(p)$$
&
\begin{center}
\begin{tikzpicture}
\sbEntree{E}
\sbBloc{sys}{$ \quad \dfrac{1}{Mp} \quad $}{E} \sbRelier[$ F(p)\quad $]{E}{sys}
\sbSortie{S}{sys} \sbRelier[$ \quad V(p)$]{sys}{S}
\end{tikzpicture}
\end{center} \\
\begin{center}
Théorème du moment dynamique (position)
\end{center}
&
\textit{Mouvement d'un solide en rotation}
%\includegraphics[height=1cm]{images/inertie_2}}
&
$$ C(t)=J\dfrac{d^2\theta(t)}{dt^2}$$
&
$$ C(p)=Jp^2\Theta(p)$$
&
\begin{center}
\begin{tikzpicture}
\sbEntree{E}
\sbBloc{sys}{$ \quad \dfrac{1}{Jp^2} \quad $}{E} \sbRelier[$ C(p)\quad $]{E}{sys}
\sbSortie{S}{sys} \sbRelier[$ \quad \Theta(p)$]{sys}{S}
\end{tikzpicture}
\end{center} \\
\begin{center}
Théorème du moment dynamique (vitesse)
\end{center}
&
\begin{center}
\includegraphics[height=1cm]{images/inertie_2}
\end{center}
&
$$ C(t)=J\dfrac{d\omega(t)}{dt}$$
&
$$ C(p)=Jp\Omega(p)$$
&
\begin{center}
\begin{tikzpicture}
\sbEntree{E}
\sbBloc{sys}{$ \quad \dfrac{1}{Jp} \quad $}{E} \sbRelier[$ C(p)\quad $]{E}{sys}
\sbSortie{S}{sys} \sbRelier[$ \quad \Omega(p)$]{sys}{S}
\end{tikzpicture} 
\end{center}
\end{tabular}
\end{center}
}

\begin{rem}
On note :
\begin{multicols}{2}
\begin{itemize}
\item $C(t)$ : le couple exprimé en $Nm$
\item $J$ : l'inertie exprimée en $kg.m^2$
\item $\theta$ : position angulaire en $rad$
\item $\omega$ : vitesse angulaire en $rad/s$
\end{itemize}
\end{multicols}
\end{rem}


%\subsection{Théorèmes de la mécaniques}
%
%\begin{minipage}[t]{.2\linewidth}
%\begin{center}
%Théorème de la résultante dynamique (position)
%
%\includegraphics[width=.3\textwidth]{images/poids}
%\end{center}
%\end{minipage} \hfill
%\begin{minipage}[t]{.2\linewidth}
%\begin{center}
%Domaine temporel
%\end{center}
%$$ F(t)=M\dfrac{d^2x(t)}{dt^2}$$
%\end{minipage} \hfill
%\begin{minipage}[t]{.2\linewidth}
%\begin{center}
%Domaine de Laplace
%\end{center}
%$$ F(p)=Mp^2X(p)$$
%\end{minipage} \hfill
%\begin{minipage}[t]{.2\linewidth}
%\begin{center}
%Schéma bloc
%
%\vspace{.5cm} 
%
%\begin{tikzpicture}
%\sbEntree{E}
%\sbBloc{sys}{$ \quad \dfrac{1}{Mp^2} \quad $}{E} \sbRelier[$ F(p)\quad $]{E}{sys}
%\sbSortie{S}{sys} \sbRelier[$ \quad X(p)$]{sys}{S}
%\end{tikzpicture} 
%\end{center}
%\end{minipage}


%\begin{minipage}[t]{.2\linewidth}
%\begin{center}
%Théorème de la résultante dynamique (vitesse)
%
%\includegraphics[width=.3\textwidth]{images/poids}
%\end{center}
%\end{minipage} \hfill
%\begin{minipage}[t]{.2\linewidth}
%\begin{center}
%Domaine temporel
%\end{center}
%$$ F(t)=M\dfrac{dv(t)}{dt}$$
%\end{minipage} \hfill
%\begin{minipage}[t]{.2\linewidth}
%\begin{center}
%Domaine de Laplace
%\end{center}
%$$ F(p)=MpV(p)$$
%\end{minipage} \hfill
%\begin{minipage}[t]{.2\linewidth}
%\begin{center}
%Schéma bloc
%
%\vspace{.5cm} 
%
%\begin{tikzpicture}
%\sbEntree{E}
%\sbBloc{sys}{$ \quad \dfrac{1}{Mp} \quad $}{E} \sbRelier[$ F(p)\quad $]{E}{sys}
%\sbSortie{S}{sys} \sbRelier[$ \quad V(p)$]{sys}{S}
%\end{tikzpicture} 
%\end{center}
%\end{minipage}




%
%
%\begin{minipage}[t]{.2\linewidth}
%\begin{center}
%Théorème du moment dynamique (position)
%
%\includegraphics[width=.75\textwidth]{images/inertie}
%\end{center}
%\end{minipage} \hfill
%\begin{minipage}[t]{.2\linewidth}
%\begin{center}
%Domaine temporel
%\end{center}
%$$ C(t)=J\dfrac{d^2\theta(t)}{dt^2}$$
%\end{minipage} \hfill
%\begin{minipage}[t]{.2\linewidth}
%\begin{center}
%Domaine de Laplace
%\end{center}
%$$ C(p)=Jp^2\Theta(p)$$
%\end{minipage} \hfill
%\begin{minipage}[t]{.2\linewidth}
%\begin{center}
%Schéma bloc
%
%\vspace{.5cm} 
%
%\begin{tikzpicture}
%\sbEntree{E}
%\sbBloc{sys}{$ \quad \dfrac{1}{Jp^2} \quad $}{E} \sbRelier[$ C(p)\quad $]{E}{sys}
%\sbSortie{S}{sys} \sbRelier[$ \quad \Theta(p)$]{sys}{S}
%\end{tikzpicture} 
%\end{center}
%\end{minipage}




%\begin{minipage}[t]{.2\linewidth}
%\begin{center}
%Théorème du moment dynamique (vitesse)
%
%\includegraphics[width=.75\textwidth]{images/inertie}
%\end{center}
%\end{minipage} \hfill
%\begin{minipage}[t]{.2\linewidth}
%\begin{center}
%Domaine temporel
%\end{center}
%$$ C(t)=J\dfrac{d\omega(t)}{dt}$$
%\end{minipage} \hfill
%\begin{minipage}[t]{.2\linewidth}
%\begin{center}
%Domaine de Laplace
%\end{center}
%$$ C(p)=Jp\Omega(p)$$
%\end{minipage} \hfill
%\begin{minipage}[t]{.2\linewidth}
%\begin{center}
%Schéma bloc
%
%\vspace{.5cm} 
%
%\begin{tikzpicture}
%\sbEntree{E}
%\sbBloc{sys}{$ \quad \dfrac{1}{Jp} \quad $}{E} \sbRelier[$ C(p)\quad $]{E}{sys}
%\sbSortie{S}{sys} \sbRelier[$ \quad \Omega(p)$]{sys}{S}
%\end{tikzpicture} 
%\end{center}
%\end{minipage}




\subsection{Lois de frottements visqueux}



\footnotesize{
\begin{center}
%\begin{tabular}{m{3cm}m{3cm}m{3cm}m{3cm}m{3cm}}
\begin{tabular}{p{3cm}p{3cm}p{3cm}p{3cm}p{3cm}}
%\multicolumn{5}{c}{\textbf{Théorèmes de la mécanique}}\\
\begin{center}
Frottements visqueux pour des solides en translation
\end{center}
&
\begin{center}
\includegraphics[width=2.5cm]{images/amortisseur}
\end{center}
&
$$ f(t)=f\dfrac{dx(t)}{dt}$$
&
$$ F(p)=fpX(p)$$
&
\begin{center}
\begin{tikzpicture}
\sbEntree{E}
\sbBloc{sys}{$ \quad \dfrac{1}{fp} \quad $}{E} \sbRelier[$ F(p)\quad $]{E}{sys}
\sbSortie{S}{sys} \sbRelier[$ \quad X(p)$]{sys}{S}
\end{tikzpicture}
\end{center} \\
\begin{center}
Frottements visqueux pour des solides en rotation
\end{center}
&
\textit{Frottements dans les paliers d'un moteur à courant continu}
%\begin{center}
%\includegraphics[height=1cm]{images/poids}
%\end{center}
&
$$ C(t)=f\dfrac{d\theta(t)}{dt}$$
&
$$ C(p)=fp\Theta(p)$$
&
\begin{center}
\begin{tikzpicture}
\sbEntree{E}
\sbBloc{sys}{$ \quad \dfrac{1}{fp} \quad $}{E} \sbRelier[$ C(p)\quad $]{E}{sys}
\sbSortie{S}{sys} \sbRelier[$ \quad \Theta(p)$]{sys}{S}
\end{tikzpicture}
\end{center} \\
\end{tabular}
\end{center}}

%\begin{minipage}[t]{.2\linewidth}
%\begin{center}
%Frottements visqueux pour des solides en translation
%
%\includegraphics[width=.6\textwidth]{images/amortisseur}
%\end{center}
%\end{minipage} \hfill
%\begin{minipage}[t]{.2\linewidth}
%\begin{center}
%Domaine temporel
%\end{center}
%$$ f(t)=f\dfrac{dx(t)}{dt}$$
%\end{minipage} \hfill
%\begin{minipage}[t]{.2\linewidth}
%\begin{center}
%Domaine de Laplace
%\end{center}
%$$ F(p)=fpX(p)$$
%\end{minipage} \hfill
%\begin{minipage}[t]{.2\linewidth}
%\begin{center}
%Schéma bloc
%
%\vspace{.5cm} 
%
%\begin{tikzpicture}
%\sbEntree{E}
%\sbBloc{sys}{$ \quad \dfrac{1}{fp} \quad $}{E} \sbRelier[$ F(p)\quad $]{E}{sys}
%\sbSortie{S}{sys} \sbRelier[$ \quad X(p)$]{sys}{S}
%\end{tikzpicture} 
%\end{center}
%\end{minipage}
%
%\begin{minipage}[t]{.2\linewidth}
%\begin{center}
%Frottements visqueux pour des solides en rotation
%
%%\includegraphics[width=.3\textwidth]{images/poids}
%\end{center}
%\end{minipage} \hfill
%\begin{minipage}[t]{.2\linewidth}
%\begin{center}
%Domaine temporel
%\end{center}
%$$ C(t)=f\dfrac{d\theta(t)}{dt}$$
%\end{minipage} \hfill
%\begin{minipage}[t]{.2\linewidth}
%\begin{center}
%Domaine de Laplace
%\end{center}
%$$ C(p)=fp\Theta(p)$$
%\end{minipage} \hfill
%\begin{minipage}[t]{.2\linewidth}
%\begin{center}
%Schéma bloc
%
%\vspace{.5cm} 
%
%\begin{tikzpicture}
%\sbEntree{E}
%\sbBloc{sys}{$ \quad \dfrac{1}{fp} \quad $}{E} \sbRelier[$ C(p)\quad $]{E}{sys}
%\sbSortie{S}{sys} \sbRelier[$ \quad \Theta(p)$]{sys}{S}
%\end{tikzpicture} 
%\end{center}
%\end{minipage}
%
%
%
%\begin{exemple}
%\begin{itemize}
%\item La loi de comportement dans un amortisseur est une loi de comportement visqueux.
%\item Dans le cas du MCC, le couple résistant $c_r(t)$ est un couple dû au frottement visqueux.
%\end{itemize}
%\end{exemple}

\subsection{Lois de comportement des ressorts}


\footnotesize{
\begin{center}
%\begin{tabular}{m{3cm}m{3cm}m{3cm}m{3cm}m{3cm}}
\begin{tabular}{p{3cm}p{3cm}p{3cm}p{3cm}p{3cm}}
%\multicolumn{5}{c}{\textbf{Théorèmes de la mécanique}}\\
\begin{center}
Ressorts en compression de raideur $k$
\end{center}
&
\begin{center}
\includegraphics[width=1cm]{images/ressort_comp}
\end{center}
&
$$ f(t)=kx(t)$$
&
$$ F(p)=kX(p)$$
&
\begin{center}
\begin{tikzpicture}
\sbEntree{E}
\sbBloc{sys}{$ \quad \dfrac{1}{k} \quad $}{E} \sbRelier[$ F(p)\quad $]{E}{sys}
\sbSortie{S}{sys} \sbRelier[$ \quad X(p)$]{sys}{S}
\end{tikzpicture}
\end{center} \\
\begin{center}
Ressorts en traction de raideur $k$
\end{center}
&
\begin{center}
\includegraphics[width=2.5cm]{images/ressort_couple}
\end{center}
&
$$ C(t)=k\theta(t)$$
&
$$ C(p)=k\Theta(p)$$
&
\begin{center}
\begin{tikzpicture}
\sbEntree{E}
\sbBloc{sys}{$ \quad \dfrac{1}{k} \quad $}{E} \sbRelier[$ C(p)\quad $]{E}{sys}
\sbSortie{S}{sys} \sbRelier[$ \quad \Theta(p)$]{sys}{S}
\end{tikzpicture}
\end{center} \\
\end{tabular}
\end{center}}
%
%\begin{minipage}[t]{.2\linewidth}
%\begin{center}
%Ressorts en compression de raideur $k$
%
%\rotatebox{90}{\includegraphics[width=.5\textwidth]{images/ressort_comp}}
%\end{center}
%\end{minipage} \hfill
%\begin{minipage}[t]{.2\linewidth}
%\begin{center}
%Domaine temporel
%\end{center}
%$$ f(t)=kx(t)$$
%\end{minipage} \hfill
%\begin{minipage}[t]{.2\linewidth}
%\begin{center}
%Domaine de Laplace
%\end{center}
%$$ F(p)=kX(p)$$
%\end{minipage} \hfill
%\begin{minipage}[t]{.2\linewidth}
%\begin{center}
%Schéma bloc
%
%\vspace{.5cm} 
%
%\begin{tikzpicture}
%\sbEntree{E}
%\sbBloc{sys}{$ \quad \dfrac{1}{k} \quad $}{E} \sbRelier[$ F(p)\quad $]{E}{sys}
%\sbSortie{S}{sys} \sbRelier[$ \quad X(p)$]{sys}{S}
%\end{tikzpicture} 
%\end{center}
%\end{minipage}
%
%\begin{minipage}[t]{.2\linewidth}
%\begin{center}
%Ressorts en traction de raideur $k$
%
%\includegraphics[width=\textwidth]{images/ressort_couple}
%\end{center}
%\end{minipage} \hfill
%\begin{minipage}[t]{.2\linewidth}
%\begin{center}
%Domaine temporel
%\end{center}
%$$ C(t)=k\theta(t)$$
%\end{minipage} \hfill
%\begin{minipage}[t]{.2\linewidth}
%\begin{center}
%Domaine de Laplace
%\end{center}
%$$ C(p)=k\Theta(p)$$
%\end{minipage} \hfill
%\begin{minipage}[t]{.2\linewidth}
%\begin{center}
%Schéma bloc
%
%\vspace{.5cm} 
%
%\begin{tikzpicture}
%\sbEntree{E}
%\sbBloc{sys}{$ \quad \dfrac{1}{k} \quad $}{E} \sbRelier[$ C(p)\quad $]{E}{sys}
%\sbSortie{S}{sys} \sbRelier[$ \quad \Theta(p)$]{sys}{S}
%\end{tikzpicture} 
%\end{center}
%\end{minipage}
%



\subsection{Systèmes électriques}


\footnotesize{
\begin{center}
%\begin{tabular}{m{3cm}m{3cm}m{3cm}m{3cm}m{3cm}}
\begin{tabular}{p{3cm}p{3cm}p{3cm}p{3cm}p{3cm}}
%\multicolumn{5}{c}{\textbf{Théorèmes de la mécanique}}\\
\begin{center}
Résistor
\end{center}
&
\begin{center}
\includegraphics[width=2cm]{images/resistance}
\end{center}
&
$$ u(t)=Ri(t)$$
&
$$ U(p)=RI(p)$$
&
\begin{center}
\begin{tikzpicture}
\sbEntree{E}
\sbBloc{sys}{$ \quad R \quad $}{E} \sbRelier[$ I(p)\quad $]{E}{sys}
\sbSortie{S}{sys} \sbRelier[$ \quad U(p)$]{sys}{S}
\end{tikzpicture}
\end{center} \\
\begin{center}
Condensateur
\end{center}
&
\begin{center}
\includegraphics[width=2cm]{images/condensateur}
\end{center}
&
$$ u(t)=\dfrac{1}{C}\int i(t) dt$$
&
$$ U(p)=\dfrac{1}{C p}I(p)$$
&
\begin{center}
\begin{tikzpicture}
\sbEntree{E}
\sbBloc{sys}{$ \quad \dfrac{1}{C p} \quad $}{E} \sbRelier[$ I(p)\quad $]{E}{sys}
\sbSortie{S}{sys} \sbRelier[$ \quad U(p)$]{sys}{S}
\end{tikzpicture}
\end{center} \\
\begin{center}
Inductance
\end{center}
&
\begin{center}
\includegraphics[width=2cm]{images/inductance}
\end{center}
&
$$ u(t)=L \dfrac{d i(t)}{dt}$$
&
$$ U(p)=LpI(p)$$
&
\begin{center}
\begin{tikzpicture}
\sbEntree{E}
\sbBloc{sys}{$ \quad Lp \quad $}{E} \sbRelier[$ I(p)\quad $]{E}{sys}
\sbSortie{S}{sys} \sbRelier[$ \quad U(p)$]{sys}{S}
\end{tikzpicture}
\end{center} \\
\end{tabular}
\end{center}}

%\begin{minipage}[t]{.2\linewidth}
%\begin{center}
%Résistance
%
%\includegraphics[width=.5\textwidth]{images/resistance}
%\end{center}
%\end{minipage} \hfill
%\begin{minipage}[t]{.2\linewidth}
%\begin{center}
%Domaine temporel
%\end{center}
%$$ u(t)=Ri(t)$$
%\end{minipage} \hfill
%\begin{minipage}[t]{.2\linewidth}
%\begin{center}
%Domaine de Laplace
%\end{center}
%$$ U(p)=RI(p)$$
%\end{minipage} \hfill
%\begin{minipage}[t]{.2\linewidth}
%\begin{center}
%Schéma bloc
%
%\vspace{.5cm} 
%
%\begin{tikzpicture}
%\sbEntree{E}
%\sbBloc{sys}{$ \quad R \quad $}{E} \sbRelier[$ I(p)\quad $]{E}{sys}
%\sbSortie{S}{sys} \sbRelier[$ \quad U(p)$]{sys}{S}
%\end{tikzpicture} 
%\end{center}
%\end{minipage}
%
%
%\begin{minipage}[t]{.2\linewidth}
%\begin{center}
%Condensateur
%
%\includegraphics[width=.5\textwidth]{images/condensateur}
%\end{center}
%\end{minipage} \hfill
%\begin{minipage}[t]{.2\linewidth}
%\begin{center}
%Domaine temporel
%\end{center}
%$$ u(t)=\dfrac{1}{C}\int i(t) dt$$
%\end{minipage} \hfill
%\begin{minipage}[t]{.2\linewidth}
%\begin{center}
%Domaine de Laplace
%\end{center}
%$$ U(p)=\dfrac{1}{C p}I(p)$$
%\end{minipage} \hfill
%\begin{minipage}[t]{.2\linewidth}
%\begin{center}
%Schéma bloc
%
%\vspace{.5cm} 
%
%\begin{tikzpicture}
%\sbEntree{E}
%\sbBloc{sys}{$ \quad \dfrac{1}{C p} \quad $}{E} \sbRelier[$ I(p)\quad $]{E}{sys}
%\sbSortie{S}{sys} \sbRelier[$ \quad U(p)$]{sys}{S}
%\end{tikzpicture} 
%\end{center}
%\end{minipage}
%
%
%
%
%\begin{minipage}[t]{.2\linewidth}
%\begin{center}
%Inductance
%
%\includegraphics[width=.5\textwidth]{images/inductance}
%\end{center}
%\end{minipage} \hfill
%\begin{minipage}[t]{.2\linewidth}
%\begin{center}
%Domaine temporel
%\end{center}
%$$ u(t)=L \dfrac{d i(t)}{dt}$$
%\end{minipage} \hfill
%\begin{minipage}[t]{.2\linewidth}
%\begin{center}
%Domaine de Laplace
%\end{center}
%$$ U(p)=LpI(p)$$
%\end{minipage} \hfill
%\begin{minipage}[t]{.2\linewidth}
%\begin{center}
%Schéma bloc
%
%\vspace{.5cm} 
%
%\begin{tikzpicture}
%\sbEntree{E}
%\sbBloc{sys}{$ \quad Lp \quad $}{E} \sbRelier[$ I(p)\quad $]{E}{sys}
%\sbSortie{S}{sys} \sbRelier[$ \quad U(p)$]{sys}{S}
%\end{tikzpicture} 
%\end{center}
%\end{minipage}

\subsection{Le moteur à courant continu en schéma bloc}
\begin{exemple}
Exprimer les équations du moteur à courant continu sous forme de schéma bloc.

\begin{minipage}[c]{.3\linewidth}
$$
c_r(t) = f_v \omega_m(t)
$$
$$
C_r(p) = f_v \Omega_m(p)
$$
\end{minipage}\hfill
\begin{minipage}[c]{.3\linewidth}
\ifthenelse{\boolean{prof}}{%
\begin{center}
\begin{tikzpicture}
\sbEntree{E}
\sbBloc{sys}{$ \quad \dfrac{1}{f_v} \quad $}{E} \sbRelier[$ C_r(p)\quad $]{E}{sys}
\sbSortie{S}{sys} \sbRelier[$ \quad \Omega_m(p)$]{sys}{S}
\end{tikzpicture} 
\end{center}
}{
\rotatebox{90}{
\begin{tabular}{p{3cm}}
\\
\end{tabular}}
}
\end{minipage} \hfill
\begin{minipage}[c]{.3\linewidth}
\end{minipage}

\begin{minipage}[c]{.3\linewidth}
$$
e(t)=K_E \cdot \omega_m(t) $$
$$
E(p)=K_E \cdot \Omega_m(p)
$$
\end{minipage}\hfill
\begin{minipage}[c]{.3\linewidth}
\ifthenelse{\boolean{prof}}{%
\begin{center}
\begin{tikzpicture}
\sbEntree{E}
\sbBloc{sys}{$ \quad \dfrac{1}{K_E} \quad $}{E} \sbRelier[$ E(p)\quad $]{E}{sys}
\sbSortie{S}{sys} \sbRelier[$ \quad\quad \Omega_m(p)$]{sys}{S}
\end{tikzpicture} 
\end{center}
}{
\rotatebox{90}{
\begin{tabular}{p{3cm}}
\\
\end{tabular}}
}
\end{minipage} \hfill
\begin{minipage}[c]{.3\linewidth}
\end{minipage}

\begin{minipage}[c]{.3\linewidth}
$$
c_m(t)=K_C \cdot i(t)
$$
$$
C_m(p)=K_C \cdot I(p)
$$
\end{minipage}\hfill
\begin{minipage}[c]{.3\linewidth}
\ifthenelse{\boolean{prof}}{%
\begin{center}
\begin{tikzpicture}
\sbEntree{E}
\sbBloc{sys}{$ \quad \dfrac{1}{K_C} \quad $}{E} 
\sbRelier[$ I(p)\quad $]{E}{sys}
\sbSortie{S}{sys} 
\sbRelier[$ \quad C_m(p)$]{sys}{S}
\end{tikzpicture} 
\end{center}
}{
\rotatebox{90}{
\begin{tabular}{p{3cm}}
\\
\end{tabular}}
}
\end{minipage} \hfill
\begin{minipage}[c]{.3\linewidth}
\end{minipage}

\begin{minipage}[c]{.3\linewidth}
$$
u_R(t)=R\cdot i(t) 
$$
$$
U_R(p)=R\cdot I(p) 
$$
\end{minipage}\hfill
\begin{minipage}[c]{.3\linewidth}
\ifthenelse{\boolean{prof}}{%
\begin{center}
\begin{tikzpicture}
\sbEntree{E}
\sbBloc{sys}{$ \quad R \quad $}{E} \sbRelier[$ I(p)\quad $]{E}{sys}
\sbSortie{S}{sys} \sbRelier[$ \quad U_r(p)$]{sys}{S}
\end{tikzpicture} 
\end{center}
}{
\rotatebox{90}{
\begin{tabular}{p{3cm}}
\\
\end{tabular}}
}
\end{minipage} \hfill
\begin{minipage}[c]{.3\linewidth}
\end{minipage}

\begin{minipage}[c]{.3\linewidth}
$$
u_L(t)=L\dfrac{di(t)}{dt}
$$
$$
U_L(p)=LpI(p)
$$
\end{minipage}\hfill
\begin{minipage}[c]{.3\linewidth}
\ifthenelse{\boolean{prof}}{%
\begin{center}
\begin{tikzpicture}
\sbEntree{E}
\sbBloc{sys}{$ \quad Lp \quad $}{E} \sbRelier[$ I(p)\quad $]{E}{sys}
\sbSortie{S}{sys} \sbRelier[$ \quad U_L(p)$]{sys}{S}
\end{tikzpicture} 
\end{center}
}{
\rotatebox{90}{
\begin{tabular}{p{3cm}}
\\
\end{tabular}}
}
\end{minipage} \hfill
\begin{minipage}[c]{.3\linewidth}
\end{minipage}

\begin{minipage}[c]{.3\linewidth}
$$
u_L(t)+u_R(t)=L\dfrac{di(t)}{dt}+R\cdot i(t)
$$
$$
U_L(p)+U_R(p)=LpI(p)+RI(p)
$$
\end{minipage}\hfill
\begin{minipage}[c]{.6\linewidth}
\ifthenelse{\boolean{prof}}{%
\begin{center}
\begin{tikzpicture}
\sbEntree{E}
\sbDecaleNoeudx[4]{E}{nE}
\sbDecaleNoeudy[-2]{nE}{nh1}
\sbDecaleNoeudy[2]{nE}{nh2}
\sbBloc{h1}{$Lp$}{nh1}
\sbBloc{h2}{$R$}{nh2}
\sbCompSum[12]{comp}{E}{+}{+}{}{}
\sbSortie[6]{S}{comp} 
\sbRelieryx{nE}{h1}
\sbRelieryx{nE}{h2}
\sbRelierxy{h1}{comp}
\sbRelierxy{h2}{comp}
\sbRelier[$I(p)$]{E}{nE}
\sbRelier[$\quad \quad  \quad  LpI(p)+RI(p)$]{comp}{S}
\end{tikzpicture}
\end{center}
}{
\rotatebox{90}{
\begin{tabular}{p{3cm}}
\\
\end{tabular}}
}
\end{minipage} \hfill
%\begin{minipage}[c]{.3\linewidth}
%\end{minipage}

\begin{minipage}[c]{.3\linewidth}
$$
u(t)-e(t)=u_L(t) + u_R(t)
$$
$$
U(p)-E(p)=U_L(p) + U_R(p) 
$$
$$
U(p)-E(p)= LpI(p)+RI(p)
$$
\end{minipage}\hfill
\begin{minipage}[c]{.3\linewidth}
\ifthenelse{\boolean{prof}}{%
\begin{center}
\begin{tikzpicture}
\sbEntree{E}
%\sbEntree{Ur}
\sbEntree{Ul}
\sbCompSum{comp}{E}{}{-}{+}{}
%\sbSortie{S}{h}
\sbRelier[$U(p)$]{E}{comp}
\sbDecaleNoeudy[4]{comp}{Ul}
\sbRelier[$E(p)$]{Ul}{comp}
\sbDecaleNoeudy[-4]{comp}{Ur}
%\sbRelier[$U_R(p)$]{Ur}{comp}
\sbBloc{h}{$ \dfrac{1}{R+Lp}$}{comp}
\sbRelier{comp}{h}
\sbSortie{S}{h}
\sbRelier[$I(p)$]{h}{S}
\end{tikzpicture}
\end{center}
}{
\rotatebox{90}{
\begin{tabular}{p{3cm}}
\\
\end{tabular}}
}
\end{minipage} \hfill
\begin{minipage}[c]{.3\linewidth}
\end{minipage}

\begin{minipage}[c]{.3\linewidth}
$$
J\dfrac{d\omega_m(t)}{dt} = c_m(t) - c_r(t)
$$
$$
Jp\Omega_m(p) = C_m(p) - C_r(p)
$$
\end{minipage}\hfill
\begin{minipage}[c]{.3\linewidth}
\ifthenelse{\boolean{prof}}{%
\begin{center}
\begin{tikzpicture}
\sbEntree{E}
%\sbEntree{Ur}
\sbEntree{Ul}
\sbCompSum{comp}{E}{}{-}{+}{}
%\sbSortie{S}{h}
\sbRelier[$C_m(p)$]{E}{comp}
\sbDecaleNoeudy[4]{comp}{Ul}
\sbRelier[$C_r(p)$]{Ul}{comp}
\sbDecaleNoeudy[-4]{comp}{Ur}
%\sbRelier[$U_R(p)$]{Ur}{comp}
\sbBloc{h}{$ \dfrac{1}{Jp}$}{comp}
\sbRelier{comp}{h}
\sbSortie{S}{h}
\sbRelier[$\quad\quad \Omega_m(p)$]{h}{S}
\end{tikzpicture}
\end{center}
}{
\rotatebox{90}{
\begin{tabular}{p{3cm}}
\\
\end{tabular}}
}
\end{minipage} \hfill
\begin{minipage}[c]{.3\linewidth}
\end{minipage}

Schéma bloc du MCC complet :

\ifthenelse{\boolean{prof}}{%
\begin{center}
\includegraphics[width=.8\textwidth]{images/mcc_bloc}
\end{center}
}{
\rotatebox{90}{
\begin{tabular}{p{10cm}}
\\
\end{tabular}}
}

\end{exemple}



\section{Manipulation des schémas blocs}




\subsection{Blocs en série}
\begin{center}
\begin{tabular}{ccc}
\begin{tikzpicture}
\sbEntree{E}
\sbBloc{h1}{$H_1(p)$}{E}
\sbBloc{h2}{$H_2(p)$}{h1} 
\sbBloc{h3}{$H_3(p)$}{h2} 
\sbSortie{S}{h3} 
\sbRelier[$E(p)\quad$]{E}{h1}
\sbRelier{h1}{h2}
\sbRelier{h2}{h3}
\sbRelier[$\quad S(p)$]{h3}{S}
\end{tikzpicture}
&
$\equiv$ 
&
\begin{tikzpicture}
\sbEntree{E}
\sbBloc{h1}{$H(p)$}{E}
\sbSortie{S}{h1} 
\sbRelier[$E(p)\quad$]{E}{h1}
\sbRelier[$\quad S(p)$]{h1}{S}
\end{tikzpicture}
\end{tabular}
\end{center}

$$
H(p) = H_1(p)\cdot H_2(p) \cdot H_3(p)
$$


\subsection{Blocs en parallèle}

\begin{minipage}[c]{.4\linewidth}
\begin{center}
\begin{tikzpicture}
\sbEntree{E}
\sbDecaleNoeudx[4]{E}{nE}
\sbDecaleNoeudy[-2]{nE}{nh1}
\sbDecaleNoeudy[2]{nE}{nh2}
\sbBloc{h1}{$H_1(p)$}{nh1}
\sbBloc{h2}{$H_2(p)$}{nh2}
\sbCompSum[12]{comp}{E}{+}{+}{}{}
\sbSortie{S}{comp} 
\sbRelieryx{nE}{h1}
\sbRelieryx{nE}{h2}
\sbRelierxy{h1}{comp}
\sbRelierxy{h2}{comp}
\sbRelier[$E(p)$]{E}{nE}
\sbRelier[$S(p)$]{comp}{S}
\end{tikzpicture}
\end{center}
\end{minipage}\hfill
\begin{minipage}[c]{.15\linewidth}
$$\equiv$$
\end{minipage}\hfill
\begin{minipage}[c]{.4\linewidth}
\begin{center}
\begin{tikzpicture}
\sbEntree{E}
\sbBloc{h1}{$H_1(p)+H_2(p)$}{E}
\sbSortie{S}{h1} 
\sbRelier[$E(p)\quad$]{E}{h1}
\sbRelier[$\quad S(p)$]{h1}{S}
\end{tikzpicture}
\end{center}
\end{minipage}


$$ S(p) = H_1(p) \cdot E(p) + H_2(p) \cdot E(p) $$
$$ H(p) = H_1(p) + H_2(p)$$

\subsection{Comparateurs en série}


\begin{minipage}[c]{.4\linewidth}
\begin{center}
\begin{tikzpicture}
\sbEntree{E}
\sbEntree{Y}
\sbEntree{Z}
\sbCompSum{comp1}{E}{}{-}{+}{}
\sbDecaleNoeudy[4]{comp1}{eY}
\sbCompSum[8]{comp2}{comp1}{}{+}{+}{}
\sbDecaleNoeudy[4]{comp2}{eZ}
\sbSortie[8]{S}{comp2} 
\sbRelier[$E(p)$]{E}{comp1}
\sbRelier[$E(p)-Y(p)$]{comp1}{comp2}
\sbRelier[$\quad \quad E(p)-Y(p)+Z(p)$]{comp2}{S}
\sbRelier[$Y(p)$]{eY}{comp1}
\sbRelier[$Z(p)$]{eZ}{comp2}
\end{tikzpicture}
\end{center}
\end{minipage}\hfill
\begin{minipage}[c]{.15\linewidth}
$$\equiv$$
\end{minipage}\hfill
\begin{minipage}[c]{.4\linewidth}
\begin{center}
\begin{tikzpicture}
\sbEntree{E}
\sbEntree{Y}
\sbEntree{Z}
\sbCompSum{comp1}{E}{}{+}{+}{}
\sbDecaleNoeudy[4]{comp1}{eY}
\sbCompSum[8]{comp2}{comp1}{}{-}{+}{}
\sbDecaleNoeudy[4]{comp2}{eZ}
\sbSortie[8]{S}{comp2} 
\sbRelier[$E(p)$]{E}{comp1}
\sbRelier[$E(p)+Z(p)$]{comp1}{comp2}
\sbRelier[$\quad \quad E(p)-Y(p)+Z(p)$]{comp2}{S}
\sbRelier[$Z(p)$]{eY}{comp1}
\sbRelier[$Y(p)$]{eZ}{comp2}
\end{tikzpicture}
\end{center}
\end{minipage}


\subsection{Déplacement des points de prélèvement}

\begin{center}
 \includegraphics[width=\textwidth]{images/prelevement1}
\end{center}

\begin{center}
 \includegraphics[width=\textwidth]{images/prelevement2}
\end{center}


\subsection{Déplacement des sommateurs}


\begin{center}
 \includegraphics[width=\textwidth]{images/sommateur1}
\end{center}

\begin{center}
 \includegraphics[width=\textwidth]{images/sommateur2}
\end{center}




\section{Fonctions de transfert des systèmes bouclés}
Il peut être parfois intéressant de modifier un diagramme fonctionnel en vue de simplifier les calculs de fonctions de transfert. Pour cela quelques règles sont à respecter.

Les manipulations suivantes n’ont pas forcément de sens physique, mais sont
symboliques.

\subsection{Systèmes en boucles fermées}
Soit la structure ci-dessous constituée en trois blocs :
\begin{center}
  \includegraphics[width=.7\textwidth]{images/CDCR}
\end{center}


\begin{resultat}
\textbf{Fonction de Transfert en Boucle Fermée}

La fonction $H(p)=\dfrac{S(p)}{E(p)} $ est appelée fonction de transfert en boucle fermée du
système. On la note \textbf{FTBF}. De manière générale, la FTBF s'exprime ainsi : 

$$
H(p)= \dfrac{\text{Chaîne directe}}{
1+\text{Chaîne directe}\times\text{Chaîne de retour}}
$$

\textbf{Attention :} le signe + du dénominateur  provient du fait que la chaîne de retour arrive avec le signe - sur le comparateur.

\end{resultat}

%\begin{demo}
%\ifthenelse{\boolean{prof}}{%
\ifprof

Montrons que dans le schéma bloc précédent, la fonction de transfert en boucle fermée peut s'écrire : 
$$
H(p)=\dfrac{K\cdot G(p)}{1+K\cdot G(p) \cdot F(p)}
$$
Recherchons la relation entre $S(p)$ et $E(p)$.

En analysant la chaîne directe :
$$ S(p)= C(p) \cdot G(p) = \varepsilon(p) \cdot K \cdot G(p)$$

En analysant la chaîne de retour : 
$$ R(p)= S(p) \cdot F(p) $$

Équation donnée par le comparateur : 
$$\varepsilon(p) = E(p) - R(p)$$
\else
\rotatebox{90}{
\begin{tabular}{p{10cm}}
\\
\end{tabular}}
\fi

%\begin{warn}
%\textbf{Attention} 
%
%Le signe \textbf{-} correspond au signe \textbf{-} du
%comparateur. Lorsqu'une perturbation intervient dans la chaîne, elle peut être
%insérée avec un comparateur portant le signe \textbf{+}. La relation est
%modifiée en conséquence. 
%\end{warn}


%\ifthenelse{\boolean{prof}}{%
\ifprof
On peut donc en déduire que :
$$\varepsilon(p) = E(p) - S(p) \cdot F(p) $$

D'où :
$$ 
S(p)= \left[ E(p) - S(p) \cdot F(p) \right] \cdot K \cdot G(p) 
$$

On a donc :
$$ 
S(p)= \dfrac{K \cdot G(p)}{1+K \cdot G(p) \cdot F(p)} E(p)
$$

Au final : 
$$ 
H(p)=\dfrac{S(p)}{E(p)}= \dfrac{K \cdot G(p)}{1+K \cdot G(p) \cdot F(p)}
$$
\else
\rotatebox{90}{
\begin{tabular}{p{10cm}}
\\
\end{tabular}}
\fi

On a donc :

\begin{minipage}[c]{.45\linewidth}
\begin{center}
\begin{tikzpicture}
\sbEntree{E}
\sbComp{comp}{E}
\sbRelier[$E(p)$]{E}{comp}
\sbBloc{B1}{$K$}{comp}
\sbBloc{B2}{$G(p)$}{B1}
\sbSortie{S}{B2}
\sbRelier{E}{comp}
\sbRelier[$\varepsilon(p)$]{comp}{B1}
\sbRelier{B1}{B2}
\sbRelier[$S(p)$]{B2}{S}

\sbDecaleNoeudy[4]{B2}{nd1}
\sbBlocr{r1}{$F(p)$}{nd1}
\sbRelieryx{B2-S}{r1}
\sbRelierxy[$R(p)$]{r1}{comp}

\end{tikzpicture}
\end{center}
\end{minipage}\hfill
\begin{minipage}[c]{.05\linewidth}
$$\equiv$$
\end{minipage}\hfill
\begin{minipage}[c]{.4\linewidth}
\begin{tikzpicture}
\sbEntree{E}
\sbBloc{B1}{$H(p)$}{E}
\sbSortie{S}{B1}
\sbRelier[$E(p) \quad$]{E}{B1}
\sbRelier[$\quad S(p)$]{B1}{S}
\end{tikzpicture}
\end{minipage}


%\end{demo}

\begin{exemple}
Calcul de la FTBF du MCC :
\begin{center}
\includegraphics[width=.8\textwidth]{images/schema_mcc_complet}
\end{center}

\ifthenelse{\boolean{prof}}{%

}{
\rotatebox{90}{
\begin{tabular}{p{10cm}}
\\
\end{tabular}}
}
\end{exemple}

\subsection{Calcul de l'erreur (ou écart)}
On reprend le système suivant :
\begin{center}
\begin{tikzpicture}
\sbEntree{E}
\sbComp{comp}{E}
\sbRelier[$E(p)$]{E}{comp}
\sbBloc{B1}{$K$}{comp}
\sbBloc{B2}{$G(p)$}{B1}
\sbSortie[8]{S}{B2}
\sbRelier{E}{comp}
\sbRelier[$\varepsilon(p)$]{comp}{B1}
\sbRelier{B1}{B2}
\sbRelier[$S(p)$]{B2}{S}

\sbDecaleNoeudy[4]{B2}{nd1}
\sbBlocr{r1}{$F(p)$}{nd1}
\sbRelieryx{B2-S}{r1}
\sbRelierxy[$R(p)$]{r1}{comp}

\end{tikzpicture}
\end{center}

\begin{resultat}
\textbf{Expression de l'erreur d'un système}

Dans le but d'exprimer l'erreur du système, il est intéressant d'exprimer la
fonction de transfert relative à l'erreur $\varepsilon(p)$. On montre que : 

$$
\varepsilon(p) = \dfrac{1}{1+K \cdot G(p) \cdot F(p)} E(p) = \dfrac{\text{1}}{
1+\text{Chaîne directe}\times\text{Chaîne de retour}} E(p)
$$
\end{resultat}
%
%\begin{exemple}
%Calcul de l'erreur du MCC  :
%
%\begin{center}
%\includegraphics[width=.8\textwidth]{images/schema_MCC}
%\end{center}
%
%\ifthenelse{\boolean{prof}}{%
%\rotatebox{90}{
%\begin{tabular}{p{10cm}}
%\\
%\end{tabular}}
%}{
%\rotatebox{90}{
%\begin{tabular}{p{10cm}}
%\\
%\end{tabular}}
%}
%\end{exemple}

\subsection{Systèmes en boucles ouvertes}

Pour l'étude du comportement du système et pour déterminer certaines des ses
performances on utilise la fonction de transfert en boucle ouverte du système
(\textbf{FTBO}). On déduit le comportement en boucle fermée (cas réel) à
partir du système en boucle ouverte (la FTBO est plus simple à manipuler). 

\begin{center}
  \includegraphics[width=.6\textwidth]{images/FTBO}
\end{center}

La boucle ouverte a donc la forme suivante : 
\begin{center}
\begin{tikzpicture}
\sbEntree{E}
\sbComp{comp}{E}
\sbRelier[$E(p)$]{E}{comp}
\sbBloc[4]{B1}{$K$}{comp}
\sbBloc[4]{B2}{$G(p)$}{B1}
\sbBloc[4]{B3}{$F(p)$}{B2}
\sbSortie{S}{B3}
\sbRelier{E}{comp}
\sbRelier[$\varepsilon(p)$]{comp}{B1}
\sbRelier{B1}{B2}
\sbRelier[$S(p)$]{B2}{B3}
\sbRelier[$\quad R(p)$]{B3}{S}
\end{tikzpicture}
\end{center}


\begin{resultat}
\textbf{Fonction de Transfert en Boucle Ouverte}

La FTBO s'écrit donc ainsi :

$$
FTBO(p) = K \cdot G(p) \cdot F(p)
$$

\textbf{Attention} à ne pas confondre la FTBO et la chaîne directe !
\end{resultat}

\subsection{Fonction de transfert boucle fermée des systèmes multi-variables}
Dans un système réel, plusieurs entrées peuvent venir modifier la sortie. Ces
entrées comprennent non seulement l'entrée principale (grandeur par rapport à
laquelle on détermine la sortie) mais aussi des entrées
supplémentaires très souvent parasites (bruit, effort résistant...).

\begin{center}
\begin{tikzpicture}
\sbEntree{E}
\sbComp{comp1}{E}
\sbRelier[$E_1(p)$]{E}{comp1}
\sbBloc{B1}{$A(p)$}{comp1}
\sbSumh{comp2}{B1}
\sbBloc{B2}{$B(p)$}{comp2}
\sbSortie[4]{S}{B2}
\sbRelier{comp1}{B1}
\sbRelier{B1}{comp2}
\sbRelier{comp2}{B2}
\sbRelier[$S(p)$]{B2}{S}

\sbDecaleNoeudy[4]{B2}{nd1}
\sbBlocr{r1}{$C(p)$}{nd1}
\sbRelieryx{B2-S}{r1}
\sbRelierxy[$R(p)$]{r1}{comp}

\sbDecaleNoeudy[-4]{comp2}{nd1}
\sbEntree{E2}
\sbRelier[$E_2(p)$]{nd1}{comp2}
\end{tikzpicture}
\end{center}


%\begin{methode}
\begin{resultat}
\textbf{Principe de superposition}

Soit un système avec deux entrées $E_1(p)$ et $E_2(p)$. Pour calculer $S(p)$ on procède de la manière suivante : 
\begin{enumerate}
\item On considère que $E_1(p)\neq 0$ et que $E_2(p)=0$.
\item On calcule alors $H_1(p)=\dfrac{S(p)}{E_1(p)}$.
\item On considère que $E_1(p) = 0$ et que $E_2(p) \neq 0$.
\item On calcule alors $H_2(p)=\dfrac{S(p)}{E_2(p)}$.
\item On conclue alors par superposition que $S(p)=E_1(p) \cdot H_1(p) + E_2(p) \cdot H_2(p)$.
\end{enumerate}

Cette méthode peut être étendue à un système à $n$ entrées. On considèrera alors à chaque étape que $n-1$ entrées sont nulles.
%\end{methode}
\end{resultat}

\begin{exemple}

\begin{minipage}[c]{.4\linewidth}

\textbf{1. On considère que $E_1(p)\neq 0$ et que $E_2(p)=0$.}

\textbf{2. On calcule alors $H_1(p)=\dfrac{S(p)}{E_1(p)}$.}


\end{minipage} \hfill
\begin{minipage}[c]{.55\linewidth}
\begin{center}
\begin{tikzpicture}
\sbEntree{E}
\sbComp{comp1}{E}
\sbRelier[$E_1(p)$]{E}{comp1}
\sbBloc{B1}{$A(p)$}{comp1}
%\sbSumh{comp2}{B1}
\sbBloc{B2}{$B(p)$}{comp2}
\sbSortie[4]{S}{B2}
\sbRelier{comp1}{B1}
%\sbRelier{B1}{comp2}
%\sbRelier{comp2}{B2}
\sbRelier{B1}{B2}
\sbRelier[$S(p)$]{B2}{S}

\sbDecaleNoeudy[4]{B2}{nd1}
\sbBlocr{r1}{$C(p)$}{nd1}
\sbRelieryx{B2-S}{r1}
\sbRelierxy[$R(p)$]{r1}{comp}

\sbDecaleNoeudy[-4]{comp2}{nd1}
\sbEntree{E2}
%\sbRelier[$E_2(p)$]{nd1}{comp2}
\end{tikzpicture}
\end{center}
\end{minipage}


\begin{minipage}[c]{.4\linewidth}

\textbf{3. On considère que $E_1(p) = 0$ et que $E_2(p) \neq 0$.}

\textbf{4. On calcule alors $H_2(p)=\dfrac{S(p)}{E_2(p)}$.}

\end{minipage} \hfill
\begin{minipage}[c]{.55\linewidth}

\begin{center}
\begin{tikzpicture}
\sbEntree{E}
\sbComp{comp1}{E}
\sbRelier[$E_2(p)\quad$]{E}{comp1}
\sbBloc{B1}{$B(p)$}{comp1}
%\sbSumh{comp2}{B1}
%\sbBloc{B2}{$B(p)$}{comp2}
\sbSortie[12]{S}{B1}
\sbRelier{comp1}{B1}
%\sbRelier{B1}{comp2}
%\sbRelier{comp2}{B2}
\sbRelier[$S(p)$]{B1}{S}


\sbDecaleNoeudy[4]{S}{nd1}
\sbBlocr{r1}{$C(p)$}{nd1}
\sbBlocr{r2}{$A(p)$}{r1}
\sbRelieryx{S}{r1}
\sbRelier{r1}{r2}
\sbRelierxy{r2}{comp}

\sbDecaleNoeudy[-4]{comp2}{nd1}
\sbEntree{E2}
%\sbRelier[$E_2(p)$]{nd1}{comp2}
\end{tikzpicture}
\end{center}
\end{minipage} 


\textbf{5. On conclue alors par superposition que $S(p)=E_1(p) \cdot H_1(p) + E_2(p) \cdot H_2(p)$.}

\ifthenelse{\boolean{prof}}{%
$$
S(p)=S_1(p)+S_2(p)=\dfrac{A(p)\cdot B(p)}{1+A(p)\cdot B(p)\cdot C(p)}\cdot
E_1(p) - \dfrac{B(p)}{1+A(p)\cdot B(p)\cdot C(p)}E_2(p)
$$
}{
\rotatebox{90}{
\begin{tabular}{p{10cm}}
\\
\end{tabular}}
}

\end{exemple}
%\begin{exemple}
%%Déterminer $\Omega
%\begin{center}
%\includegraphics[width=.8\textwidth]{images/schema_mcc_complet}
%\end{center}
%\end{exemple}

\end{document}
\section{Schéma synoptique du calcul}

Le comportement $y(t)$ d'un système sous l'effet d'un signal canonique d'entrée
$x(t)$ sera décrit dans un premier temps dans le domaine de Laplace. 

La fonction de transfert $H(p)$ du système étant connue, la transformée de
Laplace $S(p)$ du signal de sortie sera conforme au produit $H(p) \cdot E(p)$
où $E(p)$ est la transformée de Laplace de la fonction temporelle du signal
d'entrée $x(t)$.

Les signaux canoniques d'entrée $x(t)$ représentent une perturbation standard
généralement très simple (fonction échelon, rampe, sinusoïde...).

La transformée inverse donne enfin la réponse temporelle cherchée. 

\begin{center}
%  \includegraphics[width=0.9\textwidth]{eps/synoptique}
\end{center}


\end{document}
